\documentclass[a4paper,12pt]{report}

\PassOptionsToPackage{table,xcdraw}{xcolor}
\usepackage{tocloft}
\renewcommand{\cftchappresnum}{CHƯƠNG~}
\renewcommand{\cftchapaftersnum}{. }
\setlength{\cftchapnumwidth}{6em}

\usepackage[colorlinks=true, linkcolor=black, urlcolor=blue]{hyperref}
\usepackage[section]{placeins}
\usepackage{mathtools}
\usepackage{amssymb}
\usepackage{cleveref}
\usepackage{siunitx}
\usepackage{float}
\usepackage{graphicx}
\usepackage[justification=centering]{caption}
\usepackage{subcaption}
\usepackage{wallpaper}
\usepackage{nomencl}
\usepackage{fancyhdr}
\renewcommand{\chaptername}{CHƯƠNG}
\makeatletter
\renewcommand{\chaptermark}[1]{%
  \markboth{%
    \chaptername~\thechapter.\ #1%
  }{}%
}
\makeatother
\usepackage{url}
\usepackage[left=2.5cm,right=2.5cm,top=2cm,bottom=3.5cm]{geometry}

\newcommand{\UE}[1]{\renewcommand{\UE}{#1}}
\newcommand{\sujet}[1]{\renewcommand{\sujet}{#1}}
\newcommand{\titre}[1]{\renewcommand{\titre}{#1}}
\newcommand{\enseignant}[1]{\renewcommand{\enseignant}{#1}}
\newcommand{\eleves}[1]{\renewcommand{\eleves}{#1}}

\newcommand{\fairemarges}{
\makenomenclature
\pagestyle{fancy}
\fancyheadoffset{1cm}
\setlength{\headheight}{2cm}
\lhead{CS338.Q21 -- Báo cáo đồ án}
\rhead{\MakeUppercase{\leftmark}}
}

\newcommand{\fairepagedegarde}{
\begin{titlepage}
\IfFileExists{logos/background.jpg}{\ThisLRCornerWallPaper{1}{logos/background.jpg}}{}
    \centering

    {\fontsize{18}{18}\selectfont
    ĐẠI HỌC QUỐC GIA THÀNH PHỐ HỒ CHÍ MINH \\
    TRƯỜNG ĐẠI HỌC CÔNG NGHỆ THÔNG TIN
    \par}
    \vspace{1.5cm}
    {\scshape\Large \sujet \par}
    \vspace{1cm}
    \rule{\linewidth}{0.2 mm} \\[0.4 cm]
    {\huge\bfseries \titre \par}
    \rule{\linewidth}{0.2 mm} \\[1.5 cm]
    \vspace{1cm}

    \begin{minipage}{0.5\textwidth}
        \begin{flushleft} \large
        \emph{\textbf{Nhóm sinh viên:}}\\
        \eleves\\
        \end{flushleft}
    \end{minipage}
    ~
    \begin{minipage}{0.4\textwidth}
        \begin{flushright} \large
        \emph{\textbf{GVHD :}} \\
        \enseignant \\
        \end{flushright}
    \end{minipage}\\[2cm]

    \vfill
    {\large \today\par}
\end{titlepage}
}

\newcommand{\tabledematieres}{
\tableofcontents
\newpage
}

% Conditional figure include: drop in figures/<name> when available.
\newcommand{\figureorplaceholder}[3]{%
  \IfFileExists{#1}{%
    \includegraphics[width=#2]{#1}%
  }{%
    \fbox{\parbox{#2}{\centering \textit{[Placeholder]}\\\texttt{#1}\\\small #3}}%
  }%
}

\captionsetup{
    font=small,
    labelfont=bf
}

\numberwithin{figure}{chapter}
\numberwithin{table}{chapter}

\usepackage{lipsum}
\usepackage{xcolor}
\definecolor{bg}{rgb}{0.95,0.95,0.95}
\usepackage{listings}
\lstdefinestyle{bashstyle}{
    backgroundcolor=\color{bg},
    basicstyle=\ttfamily\small,
    breaklines=true,
    columns=fullflexible,
    keepspaces=true,
    language=bash,
    showstringspaces=false,
    frame=single,
    framesep=4pt,
    xleftmargin=4pt,
    xrightmargin=4pt
}
\usepackage{multirow}
\usepackage{adjustbox}
\usepackage{booktabs}
\usepackage{colortbl}
\usepackage{xurl}
\usepackage{fontspec}
\usepackage{polyglossia}
\renewcommand{\chaptername}{CHƯƠNG}
\setmainlanguage{vietnamese}
\setmainfont{texgyretermes}[
    Extension      = .otf,
    UprightFont    = *-regular,
    BoldFont       = *-bold,
    ItalicFont     = *-italic,
    BoldItalicFont = *-bolditalic
]
\usepackage{cite}
\usepackage{pgfplots}
\usepackage{pgfplotstable}
\pgfplotsset{compat=1.18}

\renewcommand{\bibname}{Tài liệu tham khảo}
\setcounter{tocdepth}{3}
\setcounter{secnumdepth}{3}

\addto\captionsvietnamese{%
  \renewcommand{\contentsname}{MỤC LỤC}
  \renewcommand{\listfigurename}{DANH SÁCH HÌNH ẢNH}
  \renewcommand{\listtablename}{DANH SÁCH BẢNG}
}

\begin{document}

\titre{Đề tài: \\
Tên tiếng Việt: Đếm đối tượng Zero-shot với mẫu chất lượng \\
Tên tiếng Anh: Zero-shot Object Counting with Good Exemplars}
\sujet{NHẬN DẠNG -- CS338.Q21 \\ BÁO CÁO ĐỒ ÁN}
\UE{ }
\enseignant{TS. Dương Việt Hằng}
\eleves{Nguyễn Công Phát -- 23521143 (Nhóm trưởng) \\
Mai Thái Bình -- 23520158 \\
Vũ Việt Cương -- 23520213}

\fairemarges
\fairepagedegarde

\tabledematieres

\listoffigures
\vspace{1cm}

\listoftables
\newpage

%---- LỜI NÓI ĐẦU ------
\chapter*{LỜI NÓI ĐẦU}
\addcontentsline{toc}{chapter}{LỜI NÓI ĐẦU}

Lời đầu tiên, nhóm chúng em xin bày tỏ lòng biết ơn sâu sắc đến cô \textbf{TS. Dương Việt Hằng}, người đã tận tình hướng dẫn, định hướng và truyền đạt những kiến thức quý báu trong môn học \textit{Nhận dạng (CS338.Q21)}. Sự góp ý chuyên môn và những buổi trao đổi của cô đã giúp nhóm hiểu rõ hơn về bài toán đếm đối tượng zero-shot, hệ thống VA-Count cũng như cách đánh giá khoa học một mô hình thị giác máy tính trên bộ dữ liệu FSC147.\\[1em]
Đồ án này hiện thực lại hệ thống VA-Count~\cite{vacount} từ bài báo gốc và mở rộng theo hai hướng nghiên cứu độc lập: (i)~tích hợp \textit{Rich Prompts}~\cite{richprompts} nhằm khai thác ngữ nghĩa chi tiết từ Multimodal LLM cho bước sinh mẫu (exemplar), và (ii)~thay thế GroundingDINO bằng \textit{YOLO-World}~\cite{yoloworld} để giải quyết nút thắt cổ chai (bottleneck) về tốc độ trong khâu phát hiện (detection). Qua đó, nhóm đã tiến hành đánh giá hệ thống trên bốn cấu hình (baseline, +Rich Prompt, +YOLO-World, +YOLO-World+Rich Prompt) nhằm phân tích ảnh hưởng của từng thành phần đến chất lượng đếm (MAE, RMSE) và thời gian suy luận (inference). Toàn bộ kết quả thực nghiệm được lấy từ log huấn luyện của chính nhóm và được dẫn nguồn rõ ràng tại Chương~\ref{ch:experiments}.\\[1em]
Chúng em cũng xin chân thành cảm ơn các bạn lớp CS338.Q21 đã cùng thảo luận và đóng góp ý kiến trong suốt quá trình thực hiện. Nhóm rất mong nhận được thêm góp ý từ cô để hoàn thiện đồ án này.\\[1em]
\begin{flushright}
\textit{Thành phố Hồ Chí Minh, \today}
\end{flushright}

%---- THÔNG TIN THÀNH VIÊN ------
\chapter*{THÔNG TIN THÀNH VIÊN}
\addcontentsline{toc}{chapter}{THÔNG TIN THÀNH VIÊN}

Bảng~\ref{tab:phan_cong} trình bày phân công nhiệm vụ trong nhóm CS338.Q21 đối với đồ án \textit{"Zero-shot Object Counting with Good Exemplars"}.

\begin{table}[H]
\centering
\renewcommand{\arraystretch}{1.25}
\begin{tabular}{|l|c|c|l|p{6cm}|}
\hline
\textbf{Họ và tên} & \textbf{MSSV} & \textbf{Đóng góp} & \textbf{Vai trò} & \textbf{Phần việc đảm nhận} \\
\hline
Nguyễn Công Phát & 23521143 & 34\% & Nhóm trưởng &
Hiện thực và thực nghiệm hai cấu hình mở rộng: VA-Count với YOLO-World và VA-Count với YOLO-World + Rich Prompt; tổ chức cấu trúc dự án và tích hợp pipeline huấn luyện/đánh giá; viết Chương~1 (Tổng quan) và Chương~5 (Kết luận). \\
\hline
Mai Thái Bình & 23520158 & 33\% & Thành viên &
Hiện thực Module Rich Prompt (Gemini 2.5 Flash) và thực nghiệm cấu hình VA-Count + Rich Prompt; chạy các thử nghiệm đánh giá MAE/RMSE; thiết kế bảng kết quả và biểu đồ phân tích; viết phần Rich Prompt trong Chương~2 và phần phân tích lỗi sai. \\
\hline
Vũ Việt Cương & 23520213 & 33\% & Thành viên &
Thực nghiệm cấu hình baseline VA-Count (GroundingDINO gốc); hiện thực module YOLO-World và kiểm thử ứng dụng demo Streamlit; thiết kế slide báo cáo; viết phần YOLO-World trong Chương~2, Chương~3 (Thực nghiệm) và Chương~4 (Demo). \\
\hline
\multicolumn{5}{|l|}{\textit{Tổng cộng: 100\%.}} \\
\hline
\end{tabular}
\caption{Phân chia công việc và mức đóng góp của các thành viên CS338.Q21}
\label{tab:phan_cong}
\end{table}

Mức độ đóng góp đã được các thành viên thảo luận, thống nhất và đánh giá khách quan dựa trên khối lượng công việc thực tế đã hoàn thành.

%==================================================================
%                          CHƯƠNG 1
%==================================================================
\chapter{TỔNG QUAN ĐỀ TÀI}
\label{ch:overview}

\section{Mô tả bài toán}
\textbf{Đếm đối tượng (Object Counting)} là bài toán cơ bản trong thị giác máy tính, yêu cầu hệ thống ước lượng số lượng các thực thể thuộc một lớp ngữ nghĩa cho trước trong một ảnh tự nhiên. Khác với các tiếp cận \textit{class-specific} (chỉ đếm một lớp được huấn luyện sẵn) hay \textit{few-shot} (yêu cầu người dùng cung cấp một vài ảnh mẫu), nhánh \textbf{Zero-shot Counting} đặt mục tiêu loại bỏ hoàn toàn dữ liệu mẫu hình ảnh được cung cấp thủ công: hệ thống chỉ nhận một ảnh đầu vào và một mô tả văn bản (\textit{text prompt}) về lớp cần đếm, sau đó tự động xác định mẫu (exemplar) và sinh \textit{bản đồ mật độ (density map)} để tính tổng số lượng đối tượng~\cite{vacount, richprompts}.

\section{Lý do chọn đề tài}
Sự bùng nổ của dữ liệu thị giác trong kỷ nguyên số đặt ra nhu cầu cấp thiết về các hệ thống phân tích hình ảnh tự động, trong đó đếm đối tượng đóng vai trò then chốt cho các ứng dụng nông nghiệp thông minh, giám sát giao thông, kiểm kê kho hàng, và phân tích hình ảnh y tế. Tuy nhiên, các tiếp cận hiện nay vẫn tồn tại nhiều hạn chế:

\begin{itemize}
    \item \textbf{Phụ thuộc vào nhãn thủ công.} Các phương pháp giám sát yêu cầu lượng lớn nhãn dạng \textit{hộp giới hạn (bounding box)} hoặc \textit{chú thích điểm (dot annotation)}, vốn tốn kém và không khả thi trên dữ liệu quy mô lớn.
    \item \textbf{Hạn chế miền đối tượng.} Các mô hình giám sát thường suy giảm đáng kể trên các lớp chưa từng xuất hiện trong huấn luyện (\textit{unseen classes}).
    \item \textbf{Tính thủ công của Few-shot.} Người dùng vẫn phải tự cung cấp ảnh mẫu, cản trở khả năng tự động hóa.
\end{itemize}
Đồ án này hướng tới ba mục tiêu chính:

\begin{enumerate}
    \item \textbf{Tự động hóa quy trình lấy mẫu:} loại bỏ hoàn toàn sự can thiệp của con người trong việc chọn mẫu (exemplar).
    \item \textbf{Học cơ chế liên kết đa phương thức:} xây dựng tương tác giữa ngữ nghĩa văn bản (text prompt) và đặc trưng (feature) hình ảnh.
    \item \textbf{Tối ưu hóa khả năng tổng quát hóa:} duy trì độ chính xác trên các lớp đối tượng mới mà mô hình chưa từng nhìn thấy trong quá trình huấn luyện.
\end{enumerate}

\section{Phát biểu bài toán}
\label{sec:problem-statement}

\paragraph{Đầu vào.}
\begin{itemize}
    \item Một ảnh màu RGB $I \in \mathbb{R}^{H \times W \times 3}$ chứa các đối tượng trong bối cảnh tự nhiên.
    \item Một văn bản (text prompt) $T$ là tên lớp tiếng Anh mô tả đối tượng cần đếm (ví dụ: \texttt{"oranges"}, \texttt{"keyboard keys"}).
    \item Tập dữ liệu chuẩn FSC147~\cite{ranjan2021fsc147} dùng để huấn luyện và đánh giá.
\end{itemize}

\paragraph{Đầu ra.}
\begin{itemize}
    \item Một bản đồ mật độ (density map) $D \in \mathbb{R}^{H \times W}$ biểu diễn phân bố không gian của các đối tượng.
    \item Tổng số lượng đối tượng $\hat{y} = \int_{H,W} D \, dx\, dy$ tương ứng với mô tả $T$.
\end{itemize}

\paragraph{Ràng buộc.}
\begin{itemize}
    \item \textbf{Dữ liệu ảnh:} ảnh RGB có thể chứa một hoặc nhiều đối tượng với mật độ rất biến thiên; mỗi lần thực thi chỉ đếm \emph{một lớp} đối tượng.
    \item \textbf{Câu lệnh:} bằng tiếng Anh; mô tả vật thể hữu hình quan sát được; không sử dụng khái niệm trừu tượng (\texttt{traffic}, \texttt{meeting}) hay danh từ không đếm được (\texttt{water}, \texttt{smoke}).
\end{itemize}

\section{Lý do lựa chọn VA-Count làm baseline}
\label{sec:why-vacount}

Trong số các phương pháp zero-shot counting hiện tại, nhóm lựa chọn VA-Count~\cite{vacount} làm nền tảng vì ba lý do chính:

\begin{enumerate}
    \item \textbf{Cơ chế tự sinh mẫu (exemplar) hoàn toàn tự động.} Khác với CounTR~\cite{countr} hay FamNet~\cite{famnet} ở chế độ \textit{few-shot} vốn yêu cầu người dùng cung cấp hộp giới hạn (bounding box) mẫu, VA-Count sử dụng GroundingDINO~\cite{groundingdino} để tự động sinh mẫu từ text prompt. Đây là tiền đề then chốt để hệ thống hoạt động theo chế độ zero-shot thực sự, phù hợp với mục tiêu loại bỏ hoàn toàn sự can thiệp thủ công.
    \item \textbf{Kiến trúc hai luồng đối lập (contrastive dual-branch).} Các phương pháp đếm trước đó như BMNet~\cite{bmnet} chỉ sử dụng mẫu dương để sinh bản đồ mật độ (density map), dẫn đến dễ nhầm lẫn giữa đối tượng cần đếm và các vật thể nền có đặc trưng thị giác (visual feature) tương tự. VA-Count giải quyết vấn đề này bằng cách bổ sung một luồng Âm (Negative Process) với cơ chế học tương phản (contrastive learning), giúp mô hình phân biệt rõ ràng giữa tiền cảnh (foreground) và nền (background). Cơ chế học tương phản này được chứng minh hiệu quả trong nhiều bài toán thị giác~\cite{simclr} và đặc biệt phù hợp với bài toán đếm nơi nhiễu nền là thách thức lớn.
    \item \textbf{Thiết kế module hóa, dễ mở rộng.} Kiến trúc VA-Count tách biệt rõ ràng giữa pha lấy mẫu (EEM) và pha sinh bản đồ mật độ (NSM). Điều này cho phép thay thế từng thành phần một cách độc lập: có thể thay GroundingDINO bằng YOLO-World trong EEM mà không cần thay đổi NSM, hoặc bổ sung module sinh prompt mà không ảnh hưởng đến bộ đếm (Counter). Tính module hóa này là lợi thế lớn so với các kiến trúc đầu-cuối (end-to-end) như ZSC~\cite{zeroshot_counting_xu}.
\end{enumerate}
Hình~\ref{fig:vis_posneq} minh họa cách hệ thống VA-Count hoạt động trên thực tế: mẫu dương (hộp giới hạn xanh lá) xác định đúng đối tượng cần đếm (strawberries), trong khi mẫu âm (hộp giới hạn đỏ) đánh dấu các vùng nền. Sự phối hợp giữa hai luồng giúp bản đồ mật độ tập trung chính xác vào đối tượng mục tiêu.

\begin{figure}[H]
    \centering
    \figureorplaceholder{figures/vis_strawberry_posneq.jpg}{0.95\textwidth}{Minh họa VA-Count trên ảnh \texttt{346.jpg}: bounding box xanh lá = positive exemplars, bounding box đỏ = negative exemplars. Overlay density map cho thấy mô hình tập trung đúng vào các quả dâu tây.}
    \caption{Minh họa hoạt động của VA-Count: positive exemplars (xanh lá) và negative exemplars (đỏ) trên ảnh \texttt{346.jpg} (lớp \emph{strawberries}). Phần overlay bên phải cho thấy bản đồ mật độ tập trung đúng vào đối tượng cần đếm.}
    \label{fig:vis_posneq}
\end{figure}

\section{Đóng góp của đồ án}
Đồ án này hiện thực lại bài báo VA-Count~\cite{vacount} từ đầu và mở rộng theo hai hướng độc lập trong cùng dòng tài liệu. Các đóng góp chính:

\begin{enumerate}
    \item \textbf{Hiện thực hệ thống VA-Count cơ sở.} Module \emph{Exemplar Enhancement Module} (EEM) sử dụng GroundingDINO~\cite{groundingdino} để sinh các hộp giới hạn (bounding box) ứng viên cho lớp đối tượng cần đếm; module \emph{Noise Suppression Module} (NSM) học hàm tương phản (contrastive loss) giữa mẫu dương và mẫu âm theo công thức trong~\cite{vacount}.
    \item \textbf{Tích hợp Rich Prompts cho khâu lấy mẫu.} Áp dụng đề xuất của~\cite{richprompts}: dùng Gemini sinh mô tả phong phú từ tên lớp đầu vào, sau đó dùng CLIP~\cite{clip} xếp hạng lại (re-rank) các đề xuất để chọn ra ảnh mẫu chất lượng cao hơn.
    \item \textbf{Tích hợp YOLO-World thay cho GroundingDINO.} Sử dụng mô hình phát hiện mở từ vựng (open-vocabulary detection) YOLO-World~\cite{yoloworld} làm bộ phát hiện (detector) thay thế trong giai đoạn EEM, cho thời gian sinh mẫu nhanh hơn nhiều lần (xem Chương~\ref{ch:experiments}).
    \item \textbf{Đánh giá đầy đủ trên FSC-147.} Huấn luyện lại từ checkpoint backbone MAE~\cite{mae} và đánh giá MAE / RMSE / thời gian suy luận (inference) trên cả bốn cấu hình (baseline, +Rich Prompt, +YOLO-World, +YOLO-World+Rich Prompt), cho phép phân tích định lượng ảnh hưởng của từng thành phần đến chất lượng đếm.
    \item \textbf{Phân tích hiện tượng bão hòa ngữ nghĩa.} Qua kết quả thực nghiệm, nhóm chỉ ra rằng GroundingDINO đã đạt gần ngưỡng bão hòa đối với thông tin text đầu vào (Rich Prompt chỉ giảm MAE $0{,}19$), trong khi YOLO-World (vốn yếu hơn về ngữ nghĩa) được hưởng lợi đáng kể từ Rich Prompt (giảm MAE $1{,}12$). Phát hiện này cung cấp cơ sở thực nghiệm cho việc lựa chọn chiến lược prompting phù hợp với từng bộ phát hiện.
\end{enumerate}
Một chi tiết cài đặt được nêu rõ để tiện đối chiếu với hai bài báo gốc khi đọc Chương~\ref{ch:method}: phiên bản Gemini sử dụng tại thời điểm cài đặt là \texttt{gemini-2.5-flash}~\cite{gemini}.

%==================================================================
%                          CHƯƠNG 2
%==================================================================
\chapter{PHƯƠNG PHÁP THỰC HIỆN}
\label{ch:method}

Chương này trình bày kiến trúc hệ thống mà nhóm hiện thực, gồm ba phần: (i) baseline \textit{VA-Count}~\cite{vacount}, (ii) mở rộng theo \textit{Rich Prompts}~\cite{richprompts} dùng LLM cho khâu sinh mẫu (exemplar), và (iii) thay GroundingDINO bằng \textit{YOLO-World}~\cite{yoloworld} ở module phát hiện (detection). Mỗi mục đều ghi rõ chi tiết cài đặt thực tế (phiên bản mô hình, hàm mất mát, siêu tham số) để bảo đảm tính tái lập.

\section{Kiến trúc tổng quan của VA-Count}
\label{sec:vacount-overview}

Sau khi nhận đầu vào (ảnh + tên lớp), hệ thống chia thành hai luồng xử lý song song:

\begin{itemize}
    \item \textbf{Luồng Dương (Positive Process):} tìm các mẫu dương đại diện cho đối tượng cần đếm. Kết quả là bản đồ mật độ dương (positive density map) $D_p$ tập trung điểm sáng tại đối tượng cần đếm.
    \item \textbf{Luồng Âm (Negative Process):} tìm các mẫu âm đại diện cho vật thể nhiễu. Kết quả là bản đồ mật độ âm (negative density map) $D_n$ tập trung điểm sáng tại các vật thể không phải mục tiêu.
\end{itemize}
Hệ thống dùng hai khối phối hợp:

\begin{enumerate}
    \item \textbf{Exemplar Enhancement Module (EEM):} tự động phát hiện và tinh chỉnh các mẫu (exemplar).
    \item \textbf{Noise Suppression Module \& Counter (NSM):} sinh bản đồ mật độ và thực hiện khử nhiễu (noise suppression).
\end{enumerate}

\begin{figure}[H]
    \centering
    \figureorplaceholder{figures/system_overview.png}{0.95\textwidth}{Tổng quan kiến trúc VA-Count: hai luồng song song positive/negative qua EEM và NSM.}
    \caption{Tổng quan kiến trúc VA-Count gồm hai luồng positive/negative qua EEM và NSM. Sơ đồ chuyển vẽ theo kiến trúc đề xuất trong VA-Count~\cite{vacount}.}
    \label{fig:system_overview}
\end{figure}

\subsection{Exemplar Enhancement Module (EEM)}
\label{subsec:eem}

\textbf{Mục tiêu:} Tự động phát hiện và tinh chỉnh các mẫu (exemplars) từ ảnh đầu vào để phục vụ cho việc đếm, loại bỏ hoàn toàn sự can thiệp thủ công trong khâu chọn mẫu.\\[1em]
EEM có ba bước:

\paragraph{Bước 1: Khởi tạo mẫu bằng GroundingDINO~\cite{groundingdino}.}
\begin{itemize}
    \item \textit{Luồng Dương:} truyền ảnh (Image) và tên lớp (Class name) của đối tượng đó (ví dụ \texttt{single dog .}). Mục tiêu là để GroundingDINO phát hiện sơ bộ các hộp giới hạn (bounding box) ứng viên dương.
    \item \textit{Luồng Âm:} truyền ảnh (Image) kèm prompt chung chứa đối tượng cần đếm. Thay vì truyền vào tên lớp (Class name) của đối tượng cần đếm, ở luồng âm ta sẽ truyền tên lớp (Ví dụ \texttt{object .}).Mục tiêu là để GroundingDINO trả về tất cả vật thể có thể có trong ảnh, bao gồm cả mục tiêu và nhiễu.
\end{itemize}

\paragraph{Bước 2: Tinh chỉnh và lọc mẫu chất lượng:} Sau khi Grounding-DINO trả về các bounding box đề xuất, hệ thống thực hiện các bước lọc để chọn ra các mẫu tốt nhất.
\begin{itemize}
    \item \textbf{Loại bỏ trùng lặp (Deduplication) (chỉ luồng Âm):} Hệ thống so sánh các bounding box tìm được ở luồng Âm với các boundinng box tìm được ở luồng Dương. Nếu bounding box ở luồng Âm nào bị chồng lấn quá nhiều (IoU cao) với các bounding box ở luồng Dương, nó sẽ bị loại bỏ. Điều này đảm bảo mẫu Âm chỉ chứa các đối tượng không phải đối tượng cần đếm.
    \item \textbf{Bộ lọc đơn đối tượng (Single-Object Filter):} Cả hai luồng đều đi qua bộ lọc này, đảm bảo các mẫu được chọn chỉ chứa duy nhất một đối tượng. Điều này được thực hiện nhờ bộ phân loại nhị phân (Binary Classifier)
\end{itemize}

\paragraph{Bước 3: Chọn top-$k$ mẫu:} Hệ thống chọn ra $k$ mẫu tốt nhất (có điểm tin cậy cao nhất) cho mỗi luồng để làm Postive Exemplars và Negative Exemplars. Mặc định: top-3 cho mỗi luồng theo điểm tin cậy (confidence score) của GroundingDINO trong baseline gốc. Phiên bản Rich Prompt mở rộng lên top-5 cho luồng Dương (Mục~\ref{sec:rich-prompts}).

\subsection{Noise Suppression Module (NSM) \& Counter}
\label{subsec:nsm}
Ngay cả khi có mẫu đúng, mô hình vẫn có thể nhầm lẫn với các đối tượng nền có màu sắc tương tự. Để giải quyết vấn đề này, NSM hoạt động như một cơ chế bảo vệ, giúp phân biệt rõ ràng giữa đối tượng cần đếm và các vật thể gây nhiễu dựa trên các mẫu đã tìm được ở EEM.\\[1em]
\textbf{Mục tiêu:} Tối đa hóa sự tương đồng giữa bản đồ mật độ dương ($D_p$) với bản đồ mật độ thực tế (Ground-truth $D_g$), đồng thời giảm thiểu sự tương đồng của bản đồ mật độ âm ($D_n$) với $D_g$. Nhờ vậy, NSM giúp hệ thống không chỉ học được ``cái gì là đối tượng'' mà còn học được ``cái gì \emph{không phải} là đối tượng''.\\[1em]
Quy trình hoạt động của NSM bao gồm hai giai đoạn chính: Tương tác đặc trưng và Tối ưu hóa hàm mất mát.

\paragraph{Pha 1: Tương tác đặc trưng và Tạo bản đồ mật độ (Feature Interaction \& Density Map Generation):} 
Hệ thống dùng một bộ đếm (counter) gồm \textit{Image Encoder}, \textit{Interaction Module} và \textit{Decoder}:
\begin{itemize}
    \item Đặc trưng (feature) ảnh đầu vào đóng vai trò \textit{Query}.
    \item Phép chiếu tuyến tính (linear projection) của đặc trưng mẫu (cả $E_p$ và $E_n$) đóng vai trò \textit{Key/Value}.
\end{itemize}
Sau cơ chế attention và Decoder, hệ thống sinh hai bản đồ mật độ:
\begin{equation}
D_p = \Gamma_{\text{decode}}(F^p_{\text{fuse}}), \qquad
D_n = \Gamma_{\text{decode}}(F^n_{\text{fuse}}).
\end{equation}
Trong đó, $F_{\text{fuse}}$ là đặc trưng sau khi đã được hợp nhất thông qua cơ chế attention.

\paragraph{Pha 2: Chiến lược Học tương phản (Contrastive Learning Strategy):}
Đây là đóng góp cốt lõi của module NSM. Thay vì chỉ tối ưu hóa sai số đếm thông thường, 
VA-Count sử dụng hàm mất mát tương phản $\mathcal{L}_C$ để “đẩy” phân bố của nhiễu ra xa khỏi 
phân bố của đối tượng mục tiêu. 
\\[1em]
Tổng hàm mất mát $\mathcal{L}_{\text{total}}$ của hệ thống là tổng của hai thành phần:
\begin{itemize}
    \item \textbf{Contrastive Loss ($\mathcal{L}_C$):} Được thiết kế dưới dạng hàm Softmax để tăng xác suất dự đoán
    đúng cho luồng dương và giảm xác suất cho luồng âm:
    \begin{equation}
    \mathcal{L}_C(D_p, D_g, D_n) =
    - \log \frac{\exp\left(\mu(D_p, D_g)\right)}
    {\exp\left(\mu(D_p, D_g)\right) + \exp\left(\mu(D_n, D_g)\right)},
    \end{equation}
    trong đó $\mu(\cdot)$ là hàm đo độ tương đồng.
    \item \textbf{Density Loss ($\mathcal{L}_D$):} Sử dụng Mean Squared Error (MSE) để đảm bảo độ chính xác của
    bản đồ mật độ dự đoán so với thực tế:
    \begin{equation}
    \mathcal{L}_D(D_p, D_g) = \frac{1}{HW} \sum_{i,j} \left(D^{(i,j)}_p - D^{(i,j)}_g\right)^2.
    \end{equation}
\end{itemize}
Hàm mục tiêu cuối cùng để huấn luyện mạng là:
\begin{equation}
\mathcal{L}_{\text{total}} = \mathcal{L}_C + \mathcal{L}_D.
\end{equation}
Nhờ cơ chế này, NSM giúp hệ thống không chỉ học được ``cái gì là đối tượng'' mà còn học được ``cái gì \emph{không phải} là đối tượng''.

\section{Cải tiến với Rich Prompts}
\label{sec:rich-prompts}

\subsection{Động lực cải tiến}

Trong baseline VA-Count, bộ phát hiện (detector) GroundingDINO~\cite{groundingdino} nhận đầu vào là tên lớp đơn giản (ví dụ \texttt{"strawberries ."}). Cách tiếp cận này bộc lộ hai hạn chế quan trọng:

\begin{itemize}
    \item \textbf{Nhập nhằng ngữ nghĩa (Semantic ambiguity).} Một tên lớp ngắn có thể khớp với nhiều loại đối tượng khác nhau: ví dụ \texttt{"key"} có thể là phím bàn phím, chìa khóa, hoặc nốt nhạc. Nghiên cứu về prompt engineering trong mô hình thị giác-ngôn ngữ (vision-language)~\cite{coop} chỉ ra rằng các mô tả chi tiết (rich description) giúp giảm đáng kể sự nhập nhằng này trong không gian nhúng (embedding space).
    \item \textbf{Thiếu thông tin phân biệt tiền cảnh/nền.} Khi ảnh có nền phức tạp (nhiều vật thể cùng màu, cùng hình dạng), một tên lớp đơn không cung cấp đủ đặc trưng để GroundingDINO phân biệt đối tượng cần đếm với các vật thể gây nhiễu. Đây là vấn đề đã được nhận diện trong các nghiên cứu về phát hiện mở từ vựng (open-vocabulary detection)~\cite{ovdet_survey}.
\end{itemize}
Bài báo Rich Prompts~\cite{richprompts} giải quyết cả hai vấn đề bằng cách sử dụng một Multimodal LLM để sinh prompt mô tả chi tiết đặc trưng thị giác (visual feature) cốt lõi của đối tượng. Nhóm hiện thực ý tưởng này với \textbf{Gemini~2.5~Flash} ~\cite{gemini}, là phiên bản có sẵn ở thời điểm cài đặt. Nhóm tận dụng khả năng hiểu ngữ cảnh của mô hình ngôn ngữ lớn Gemini 2.5 Flash để sinh ra các prompt (câu nhắc) chi tiết hơn, phục vụ việc trích xuất \texttt{positive boxes} và \texttt{negative boxes}. Quy trình tổng quan được minh họa trong hình ~\ref{fig:enhanced_prompt}.

\subsection{Module Enhanced Prompt}
\label{subsec:enhanced-prompt}

Module nhận ba loại đầu vào:
\begin{itemize}
    \item \textbf{Tên lớp (Class name):} tên lớp đối tượng.
    \item \textbf{Ảnh (Image):} ảnh gốc.
    \item \textbf{Prompt hệ thống (System prompt):} gồm hai mẫu, đó là \textit{Positive Prompt Template} (mô tả đặc trưng thị giác cốt lõi) và \textit{Negative Prompt Template} (mô tả nền/nhiễu).
\end{itemize}

\begin{figure}[H]
    \centering
    \figureorplaceholder{figures/enhanced_prompt_flow.png}{0.85\textwidth}{Sơ đồ pipeline của Module Enhanced Prompt: Image+ClassName+SystemPrompt → Gemini → enhanced positive/negative prompts.}
    \caption{Sơ đồ pipeline của Module Enhanced Prompt. Sơ đồ chuyển vẽ theo thiết kế trong Rich Prompts~\cite{richprompts}.}
    \label{fig:enhanced_prompt}
\end{figure}
Chi tiết các prompt mẫu được thiết kế để định hướng mô hình tập trung vào
đặc trưng thị giác cốt lõi:

\begin{lstlisting}[style=bashstyle,caption={Mẫu Positive Prompt (Feature extraction)},label=lst:pos_prompt]
Look at the image and provide the visual definition
of a single '{singular_name}'.

Task: Describe the intrinsic physical appearance of
just ONE instance, as if it were isolated.

STRICT RULES:
1. Subject: Describe ONE item only.
   Ignore the background or other identical items.
2. Format: Start with 'single {singular_name}'.
   Use dot-separated phrases.
3. Content: Focus on Shape, Color, Material.

Example for 'keyboard key':
BAD:  keyboard key . rows of buttons . full layout .
GOOD: single keyboard key . square shape .
      black plastic material . white printed letter .

Your output for '{singular_name}':
\end{lstlisting}

\begin{lstlisting}[style=bashstyle,caption={Mẫu Negative Prompt (Loại bỏ nhiễu)},label=lst:neg_prompt]
Analyze the image and identify visual elements that
are NOT the object '{class_name}'.

Describe the background, environment, and potential
distractors.

Focus strictly on:
1. Background materials.
2. Environmental conditions.
3. Other objects that are NOT '{class_name}'.

Output format rules:
- Do NOT describe the '{class_name}' itself.
- Separate phrases with ' . '.
- Do NOT write full sentences.
- End with a dot.

Example input:  Image of a cat on a sofa.
Example output: beige sofa texture . blurry living
               room . shadows . fabric patterns .

Your output for this image:
\end{lstlisting}


\begin{figure}[H]
    \centering
    \begin{subfigure}{0.42\textwidth}
        \centering
        \figureorplaceholder{figures/positive_keyboard.png}{\linewidth}{Kết quả top-5 positive boxes cho lớp \texttt{keyboard keys} sau Re-ranking bằng CLIP ViT-B/32.}
        \caption{Class name: keyboard keys}
    \end{subfigure}\hfill
    \begin{subfigure}{0.42\textwidth}
        \centering
        \figureorplaceholder{figures/positive_grapes.png}{\linewidth}{Kết quả top-5 positive boxes cho lớp \texttt{grapes}.}
        \caption{Class name: grapes}
    \end{subfigure}
    \caption{Minh họa một số kết quả thực nghiệm về việc trích xuất \texttt{positive boxes} thành công nhờ sử dụng các prompt cải tiến trên.}
    \label{fig:positive_boxes}
\end{figure}

\subsection{Kiến trúc cải tiến}
Mô hình đề xuất bao gồm các thay đổi chiến lược trong việc trích xuất các hộp bao (bounding boxes) nhằm nâng cao độ chính xác của đầu vào cho Image Encoder. Sơ đồ kiến trúc tổng quan được trình bày trong hình~\ref{fig:system_richprompt}

\begin{figure}[H]
    \centering
    \figureorplaceholder{figures/system_richprompt.png}{0.95\textwidth}{Kiến trúc VA-Count + Rich Prompt: Gemini sinh prompt → GroundingDINO → CLIP re-rank → NSM.}
    \caption{Kiến trúc hệ thống sau khi tích hợp Rich Prompt và CLIP Re-ranking.}
    \label{fig:system_richprompt}
\end{figure}
Các cải tiến chính được chia thành hai luồng xử lý độc lập, cụ thể như sau:\\[1em]
\textbf{Luồng Dương - Cải tiến trích xuất Positive Boxes:}
\begin{itemize}
    \item \textit{Câu truy vấn đầu vào (Input Query).} Dùng \emph{positive enhanced prompt} thay vì tên lớp đơn thuần. Prompt chi tiết giúp GroundingDINO nhận diện đúng đặc trưng thị giác cốt lõi của đối tượng (ví dụ: mô tả ``square shape, black plastic material, white printed letter'' cho phím bàn phím thay vì chỉ ``keyboard key''), từ đó tăng khả năng tách biệt đối tượng mục tiêu khỏi các vật thể gây nhiễu.
    \item \textit{Xếp hạng lại (Re-ranking).} Tích hợp \emph{CLIP ViT-B/32}~\cite{clip} để đánh giá lại độ tương đồng ngữ nghĩa (semantic similarity) giữa tên lớp gốc và từng patch ảnh được cắt từ hộp giới hạn mà GroundingDINO đề xuất. Khi sử dụng enhanced prompt, GroundingDINO trở nên ``nhạy'' hơn và có xu hướng trích xuất cả các vật thể có đặc tính thị giác tương tự nhưng không phải mục tiêu (ví dụ: nút áo khi tìm phím bàn phím). Do đó, việc lấy trực tiếp top-$k$ theo logit của GroundingDINO có thể dẫn đến sai sót. CLIP đóng vai trò bộ lọc ngữ nghĩa cấp hai (second-level semantic filter), đảm bảo các patch được chọn thực sự khớp với khái niệm đối tượng (tên lớp gốc) chứ không chỉ khớp với mô tả visual chi tiết trong enhanced prompt.
    \item \textit{Chiến lược chọn mẫu.} Thay vì top-3 như VA-Count gốc, lấy \textbf{top-5} patch có điểm CLIP cao nhất. Việc mở rộng từ top-3 lên top-5 là một thử nghiệm thực nghiệm nhằm đánh giá xem việc cung cấp thêm ngữ cảnh hình ảnh (image context) cho Image Encoder có mang lại sự cải thiện đáng kể về hiệu suất đếm hay không. Trực giác là: với enhanced prompt, GroundingDINO sinh ra nhiều ứng viên chất lượng hơn, nên mở rộng $k$ giúp tận dụng được nguồn mẫu phong phú này mà không ảnh hưởng đến nhiễu (nhờ CLIP đã lọc ở bước trước).
\end{itemize}
\textbf{Luồng Âm - Cải tiến trích xuất Negative Boxes:}
\begin{itemize}
    \item \textit{Câu truy vấn đầu vào.} Thử nghiệm \emph{negative enhanced prompt} thay query mặc định \texttt{"object ."}.
    \item \textit{Loại bỏ trùng lặp (Deduplication).} Vẫn giữ tên lớp gốc (không phải enhanced prompt) làm tham chiếu cho bước loại bỏ trùng lặp. Lý do: nếu dùng enhanced prompt (vốn mô tả chi tiết đặc trưng thị giác của đối tượng) làm tham chiếu cho thuật toán NMS, nhiều hộp nền (background box) cần thiết cho việc huấn luyện luồng Âm sẽ bị loại nhầm do có điểm tương đồng cao với mô tả chi tiết. Giữ tên lớp gốc đảm bảo chỉ loại bỏ các hộp thực sự trùng với đối tượng mục tiêu.
\end{itemize}



\paragraph{Lý giải lựa chọn kỹ thuật.}
\begin{enumerate}
    \item \textbf{CLIP Re-ranking là cần thiết.} Khi dùng enhanced prompt (ví dụ: mô tả ``square shape, black plastic material, white printed letter'' cho phím bàn phím), GroundingDINO trở nên ``nhạy'' hơn và trích xuất cả các vật có đặc tính thị giác tương tự nhưng không phải mục tiêu. Việc lấy trực tiếp top-3 theo logit từ GroundingDINO có thể dẫn đến sai sót vì phân bố điểm tin cậy đã bị dịch chuyển bởi enhanced prompt. CLIP đóng vai trò bộ lọc ngữ nghĩa cấp hai (second-level semantic filter), đảm bảo các patch được chọn thực sự khớp với khái niệm đối tượng gốc (tên lớp).
    \item \textbf{Loại bỏ trùng lặp theo tên lớp gốc.} Nếu dùng enhanced prompt (vốn sinh quá nhiều positive match) làm tham chiếu cho NMS, thuật toán sẽ triệt tiêu nhầm nhiều vùng nền cần thiết cho huấn luyện luồng Âm. Giữ tên lớp gốc làm tham chiếu giúp chỉ loại bỏ các hộp thực sự trùng với đối tượng mục tiêu.
    \item \textbf{Top-3 $\to$ Top-5.} Đây là thử nghiệm thực nghiệm nhằm đánh giá xem việc cung cấp thêm ngữ cảnh hình ảnh cho Image Encoder có mang lại cải thiện đáng kể hay không. Động lực chính là: với enhanced prompt, GroundingDINO sinh nhiều ứng viên chất lượng hơn, do đó mở rộng $k$ lên 5 giúp tận dụng nguồn mẫu phong phú này. Nhiễu phát sinh từ việc tăng số lượng ứng viên được kiểm soát bởi CLIP Re-ranking ở bước trước (chọn top-5 theo điểm CLIP chứ không theo logit GroundingDINO), nên chất lượng mẫu cuối cùng vẫn được đảm bảo.
\end{enumerate}

\section{Cải tiến với YOLO-World}
\label{sec:yolo-world}

\subsection{Động lực thay thế GroundingDINO}

Trong baseline VA-Count, GroundingDINO~\cite{groundingdino} là nút thắt cổ chai (bottleneck) về tốc độ. Với kiến trúc Transformer kết hợp hợp nhất đa phương thức (cross-modal fusion), mô hình cần tính toán tầng hợp nhất thị giác-ngôn ngữ (Vision-Language Fusion) cho mỗi frame trong thời gian thực, dẫn đến độ trễ (latency) cao và tiêu tốn tài nguyên đáng kể. Cụ thể, trong pipeline VA-Count, bước trích xuất mẫu chiếm tới $\sim$80\% tổng thời gian xử lý mỗi ảnh.
Vấn đề này đặc biệt nghiêm trọng khi triển khai thực tế:
\begin{itemize}
    \item \textbf{Thời gian chuẩn bị dữ liệu.} Trên  toàn bộ 6\,135 ảnh FSC-147, GroundingDINO tốn $\approx$13 giờ cho cả hai luồng (positive + negative). Con số này tăng lên 25 giờ khi bật Rich Prompt. Đây là rào cản lớn cho việc thử nghiệm nhanh các cấu hình khác nhau.
    \item \textbf{Ứng dụng tương tác.} Khi chạy demo, người dùng cần chờ $\sim$1,2 giây/ảnh chỉ cho bước phát hiện (tương đương $\sim$80\% tổng pipeline 1,47 s; xem Bảng~\ref{tab:demo_time}), chưa kể bước sinh bản đồ mật độ. Yêu cầu phản hồi nhanh trong ứng dụng thực tế đòi hỏi thời gian xử lý dưới 1 giây tổng.
\end{itemize}

\subsection{Lựa chọn YOLO-World}

Nhóm thay bộ phát hiện (detector) bằng \emph{YOLO-World}~\cite{yoloworld}, một mô hình phát hiện mở từ vựng (open-vocabulary) nhẹ. Theo bài báo gốc, YOLO-World nhanh hơn GroundingDINO khoảng 20$\times$ nhờ chiến lược \textit{Prompt-then-Detect} và kiến trúc RepPAN (Re-parameterizable Path Aggregation Network) - con số này được đo trên tập chuẩn LVIS minival trong điều kiện phòng thí nghiệm. Trong pipeline demo của nhóm, mức tăng tốc thực tế là $\approx$2,5$\times$ so với VA-Count gốc (xem Bảng~\ref{tab:demo_time}), thấp hơn do pipeline tích hợp thêm các bước xử lý (CLIP Re-ranking, NSM, bản đồ mật độ) vốn không có trong điều kiện đo chuẩn của bài báo.\\[1em]
Cốt lõi của chiến lược \textit{Prompt-then-Detect} nằm ở việc mã hóa ngoại tuyến (offline encode) các text prompt thành tham số của mô hình trước khi chạy suy luận (inference). Cụ thể, trong quá trình chuẩn bị, bộ mã hóa ngôn ngữ (language encoder) của YOLO-World xử lý tất cả các prompt văn bản và nhúng (embed) chúng thành các vector đặc trưng cố định. Các vector này sau đó được tích hợp trực tiếp vào các tầng RepPAN của mạng phát hiện. Nhờ vậy, trong quá trình suy luận, mô hình hoàn toàn không cần gọi lại bộ mã hóa ngôn ngữ, loại bỏ gánh nặng tính toán đáng kể và đưa tốc độ xử lý về mức tương đương mạng CNN truyền thống.\\[1em]
Kiến trúc RepPAN (Re-parameterizable Path Aggregation Network) đóng vai trò then chốt trong việc hợp nhất đặc trưng đa tỉ lệ (multi-scale feature) từ backbone với thông tin ngôn ngữ đã được mã hóa. Thay vì thực hiện tương tác thị giác-ngôn ngữ (vision-language interaction) tại mỗi bước suy luận như GroundingDINO, RepPAN ``đóng gói'' thông tin ngôn ngữ thành tham số mạng cố định thông qua kỹ thuật tái tham số hóa (re-parameterization), cho phép quá trình phát hiện diễn ra hoàn toàn trên nhánh thị giác mà không cần truy vấn lại nhánh ngôn ngữ.\\[1em]
Hình~\ref{fig:vis_dense} minh họa một trường hợp đếm thực tế với mật độ dày đặc (175 đối tượng): mẫu âm (đỏ, góc trên trái; 1 trong 3 hộp âm được hiển thị) đánh dấu vùng nền, mẫu dương (xanh lá) chọn đúng marker pen. Dù mô hình đếm thiếu đáng kể (dự đoán 124,89), ví dụ này cho thấy rõ vai trò của NSM trong việc phân biệt tiền cảnh/nền.

\begin{figure}[H]
    \centering
    \figureorplaceholder{figures/vis_dense_negexemplar.jpg}{0.95\textwidth}{Minh họa VA-Count trên ảnh dày đặc \texttt{6822.jpg}: positive exemplars (xanh lá) và negative exemplar (đỏ). GT=175, Predicted=124.89. Trường hợp đếm thiếu do mật độ quá dày.}
    \caption{Trường hợp mật độ dày đặc (\texttt{6822.jpg}, lớp \emph{marker pens}): mẫu dương (xanh lá) và mẫu âm (đỏ, góc trên trái; chỉ 1/3 hộp âm hiển thị). GT $= 175$, dự đoán $\approx 124{,}89$. Đây là ví dụ điển hình về lỗi đếm thiếu do vật thể quá nhỏ và sát nhau.}
    \label{fig:vis_dense}
\end{figure}

\begin{figure}[H]
    \centering
    \figureorplaceholder{figures/yolo_speed_curve.png}{0.65\textwidth}{Speed-and-Accuracy Curve so sánh YOLO-World với GroundingDINO/GLIP trên tập LVIS minival (trích từ~\cite{yoloworld}).}
    \caption{Speed-and-Accuracy Curve: YOLO-World so với GroundingDINO/GLIP trên LVIS minival. Đồ thị trích từ bài báo gốc YOLO-World~\cite{yoloworld}.}
    \label{fig:yolo_speed}
\end{figure}
YOLO-World đạt 35 mAP với 50--80 FPS trên LVIS minival, so với 27 mAP và <5 FPS của GroundingDINO. Khi tích hợp YOLO-World vào VA-Count, hai luồng EEM được thay GroundingDINO tương ứng (\texttt{yolo\_pos\_withPrompt.py}, \texttt{yolo\_pos\_withoutPrompt.py}, \texttt{yolo\_neg.py}), trong khi NSM giữ nguyên.

\begin{figure}[H]
    \centering
    \figureorplaceholder{figures/system_yolo.png}{0.95\textwidth}{Kiến trúc VA-Count với YOLO-World thay thế GroundingDINO. Nhánh prompt enhancement vẫn được dùng tùy chọn.}
    \caption{Kiến trúc hệ thống sau khi thay GroundingDINO bằng YOLO-World ở module phát hiện. Sơ đồ chuyển vẽ theo thiết kế của VA-Count~\cite{vacount} với phần detection được thay theo YOLO-World~\cite{yoloworld}.}
    \label{fig:system_yolo}
\end{figure}

%==================================================================
%                          CHƯƠNG 3
%==================================================================
\chapter{THỰC NGHIỆM VÀ KẾT QUẢ}
\label{ch:experiments}

\section{Bộ dữ liệu FSC147}

FSC147~\cite{ranjan2021fsc147} là bộ dữ liệu chuẩn cho \textit{đếm bất kể lớp / zero-shot (class-agnostic / zero-shot counting)}: 6\,135 ảnh, 147 lớp, mỗi ảnh được chú thích các \textit{chú thích điểm (dot annotation)}. Bộ dữ liệu nổi bật với:
\begin{itemize}
    \item Các tập con không chồng lấn (train/val/test theo lớp).
    \item Độ thách thức cao do biến thiên lớn về mật độ và kích thước đối tượng.
    \item Nhãn lớp + chú thích dạng điểm cho phép đánh giá zero-shot thông qua text prompt.
\end{itemize}
Nhóm áp dụng chiến lược chia 6:2:2 (train:val:test) theo split chính thức của FSC-147~\cite{ranjan2021fsc147}.

\begin{figure}[H]
    \centering
    \figureorplaceholder{figures/fsc147_samples.png}{0.85\textwidth}{Một số ảnh minh họa từ FSC147: từ các lớp dày đặc (oranges, beans, keyboard keys) đến các lớp thưa.}
    \caption{Một vài ảnh minh họa từ bộ dữ liệu FSC-147~\cite{ranjan2021fsc147}.}
    \label{fig:fsc147}
\end{figure}
Hình~\ref{fig:vis_good} và~\ref{fig:vis_peas} minh họa kết quả đếm định tính trên hai ảnh từ tập test. Hệ thống cho thấy khả năng đếm chính xác trên nhiều loại đối tượng đa dạng, với bản đồ mật độ tập trung đúng vào vị trí các đối tượng.

\begin{figure}[H]
    \centering
    \figureorplaceholder{figures/vis_good_prediction.png}{0.95\textwidth}{Kết quả tốt: ảnh 4910.jpg (chocolates), GT=14, Predicted=14.40, sai số 2.86\%.}
    \caption{Kết quả đếm tốt trên ảnh \texttt{4910.jpg} (lớp \emph{chocolates}): 3 mẫu dương được chọn đúng, bản đồ mật độ phân bố chính xác. GT $= 14$, dự đoán $= 14{,}40$ (sai số $2{,}86\%$).}
    \label{fig:vis_good}
\end{figure}

\begin{figure}[H]
    \centering
    \figureorplaceholder{figures/vis_peas_result.png}{0.95\textwidth}{Kết quả đếm trên ảnh 5511.jpg (peas), GT=18, Predicted=19.48, sai số 8.22\%.}
    \caption{Kết quả đếm trên ảnh \texttt{5511.jpg} (lớp \emph{peas}): dù đối tượng nhỏ và xếp dày, hệ thống vẫn duy trì sai số thấp. GT $= 18$, dự đoán $= 19{,}48$ (sai số $8{,}22\%$).}
    \label{fig:vis_peas}
\end{figure}

\section{Độ đo đánh giá}

Hệ thống được đánh giá bằng hai độ đo (metric) chuẩn cho bài toán đếm: MAE và RMSE.

\subsection{Mean Absolute Error (MAE)}
MAE đo trung bình giá trị sai số tuyệt đối:
\begin{equation}
\mathrm{MAE} = \frac{1}{N}\sum_{i=1}^{N}\left|y_i - \hat{y}_i\right|,
\end{equation}
trong đó $N$ là số ảnh, $y_i$ là số đối tượng thực tế (ground truth), $\hat{y}_i$ là số đếm dự đoán.

\subsection{Root Mean Squared Error (RMSE)}
RMSE phạt nặng các sai số lớn nhờ bình phương trước khi lấy căn:
\begin{equation}
\mathrm{RMSE} = \sqrt{\frac{1}{N}\sum_{i=1}^{N}(y_i - \hat{y}_i)^2}.
\end{equation}
Sử dụng đồng thời cả hai cho phép đánh giá toàn diện: MAE phản ánh độ chính xác trung bình; RMSE phản ánh độ ổn định trên các trường hợp khó.

\section{Thiết lập thực nghiệm}
\label{sec:setup}

Các con số trong các bảng dưới đây được thu thập trên cấu hình:
\begin{itemize}
    \item \textbf{Huấn luyện và đánh giá:} server NVIDIA H200, 24 nhân CPU, 192\,GB RAM, 3\,TB NVMe.
    \item \textbf{Demo:} laptop Intel Core i7-13620H + RTX~4060.
    \item \textbf{Trích xuất mẫu (exemplar extraction):} thực hiện trên máy demo (RTX~4060), vì đây là bước tiền xử lý độc lập và không phụ thuộc môi trường huấn luyện. Thời gian tương ứng được liệt kê trong Bảng~\ref{tab:extract_time}.
\end{itemize}

\begin{table}[H]
\centering
\renewcommand{\arraystretch}{1.2}
\begin{tabular}{|l|c|}
\hline
\rowcolor[HTML]{EFEFEF}
\textbf{Tham số} & \textbf{Giá trị} \\
\hline
Bộ tối ưu (Optimizer) & AdamW \\
\hline
Tốc độ học cơ sở (Base Learning Rate) & $1{\times}10^{-5}$ \\
\hline
Suy giảm trọng số (Weight Decay) & $0.05$ \\
\hline
Lịch tốc độ học (LR Schedule) & Cosine Annealing \\
\hline
Tổng số epoch & $500$ \\
\hline
Kích thước batch (per GPU) & $8$ \\
\hline
Tích lũy gradient (Gradient Accumulation) & $16$ bước \\
\hline
\end{tabular}
\caption{Siêu tham số (hyperparameter) huấn luyện được nhóm sử dụng, theo cấu hình khuyến nghị trong VA-Count~\cite{vacount}.}
\label{tab:hyperparams}
\end{table}

\paragraph{Lưu ý.}
CLI mặc định trong \texttt{FSC\_train.py} có giá trị khác
(epochs $1000$, lr $5{\times}10^{-5}$).
Bảng~\ref{tab:hyperparams} liệt kê tổ hợp đã được nhóm sử dụng cho các
kết quả báo cáo, lưu trong log \texttt{wandb} ở thư mục
\texttt{experiments/exp\{2,3,4,5\}/wandb/}.\footnote{Thực nghiệm exp1 là lần chạy thử ban đầu chưa có logging đầy đủ (không ghi \texttt{wandb-summary.json}) và không được đưa vào phân tích chính thức.}

\section{Kết quả thực nghiệm}
\label{sec:results}

Các kết quả ở Bảng~\ref{tab:extract_time}, Bảng~\ref{tab:counting_perf},
Bảng~\ref{tab:demo_time} được tổng hợp từ các log \texttt{wandb} của nhóm
trong \mbox{\texttt{experiments/exp\{2,3,4,5\}/wandb/}}
và đã được đối chiếu lại với các tệp \texttt{wandb-summary.json} còn lưu được.

\subsection{Thời gian trích xuất mẫu (exemplar)}

\begin{table}[H]
\centering
\renewcommand{\arraystretch}{1.2}
\begin{tabular}{|l|c|c|}
\hline
\rowcolor[HTML]{EFEFEF}
\textbf{Phương pháp} & \textbf{Positive (giờ)} & \textbf{Negative (giờ)} \\
\hline
Grounding DINO (Baseline) & $\approx 3$ & $\approx 10$ \\
\hline
Grounding DINO + Rich Prompt & $\approx 10$ & $\approx 15$ \\
\hline
YOLO-World & $\approx 0.5$ & $\approx 2$ \\
\hline
YOLO-World + Rich Prompt & $\approx 1$ & $\approx 3$ \\
\hline
\end{tabular}
\caption{Thời gian trích xuất mẫu trên toàn bộ 6\,135 ảnh FSC147 (đo bởi nhóm trên máy demo RTX~4060; bước trích xuất mẫu là tiền xử lý độc lập với môi trường huấn luyện H200).}
\label{tab:extract_time}
\end{table}

YOLO-World hoàn thành trích xuất hộp dương (positive box) chỉ trong $\approx 0{.}5$ giờ, nhanh gấp 6 lần GroundingDINO ($\approx 3$ giờ). Khi tích hợp Rich Prompt, ưu thế càng rõ: GroundingDINO + Rich Prompt mất tới 25 giờ tổng, trong khi YOLO-World + Rich Prompt chỉ tốn 4 giờ. Đây là một khác biệt mang tính quyết định cho việc chuẩn bị dữ liệu trên bộ dữ liệu lớn.

\subsection{Hiệu suất đếm trên FSC-147}

\begin{table}[H]
\centering
\renewcommand{\arraystretch}{1.2}
\begin{tabular}{|l|c|c|}
\hline
\rowcolor[HTML]{EFEFEF}
\textbf{Mô hình} & \textbf{MAE} $\downarrow$ & \textbf{RMSE} $\downarrow$ \\
\hline
VA-Count~\cite{vacount} (gốc) & 17,88 & 129,31 \\
\hline
VA-Count & 17,99 & \textbf{129,39} \\
\hline
VA-Count + Rich Prompt & \textbf{17,80} & 129,69 \\
\hline
VA-Count với YOLO-World & 19,03 & 131,55 \\
\hline
VA-Count với YOLO-World + Rich Prompt & 17,91 & 130,98 \\
\hline
\end{tabular}
\caption{So sánh hiệu suất đếm trên FSC-147 (test split). Dòng ``(gốc)'' lấy từ bài báo VA-Count~\cite{vacount} để đối chiếu; các dòng còn lại là kết quả của nhóm, tổng hợp từ log \texttt{wandb} trong thư mục \mbox{\texttt{experiments/exp\{2,3,4,5\}/wandb/}}.}
\label{tab:counting_perf}
\end{table}

\paragraph{Phân tích kết quả.}

Nhìn tổng thể, bốn cấu hình cho kết quả MAE trong khoảng hẹp $17{,}80$--$19{,}03$, với RMSE dao động $129{,}39$--$131{,}55$. Mức chênh lệch này tương đối nhỏ so với giá trị tuyệt đối, gợi ý rằng \emph{chất lượng đếm cuối cùng bị giới hạn chủ yếu bởi module bộ đếm (Counter) dựa trên bản đồ mật độ}, chứ không phải bởi nguồn mẫu (exemplar source). Cụ thể, cả bốn cấu hình đều dùng chung kiến trúc Counter MAE-based (Image Encoder + Interaction Module + Decoder) và chỉ khác nhau ở khâu trích xuất mẫu (EEM). Điều này đặt ra một giới hạn trên (upper bound) cho mức cải thiện mà bất kỳ thay đổi nào ở phía bộ phát hiện có thể mang lại.

\begin{itemize}
    \item \textbf{Rich Prompt nâng nhẹ chất lượng đếm:} MAE giảm từ 17,99 xuống 17,80 trên VA-Count gốc ($\Delta = 0{,}19$). Mức cải thiện khiêm tốn này cho thấy GroundingDINO vốn đã có năng lực hiểu ngữ nghĩa mạnh nhờ kiến trúc hợp nhất đa phương thức (cross-modal fusion), và đã đạt gần \emph{ngưỡng bão hòa ngữ nghĩa} (semantic saturation) đối với thông tin đầu vào dạng văn bản. Nói cách khác, GroundingDINO đã khai thác gần hết lượng thông tin ngữ nghĩa có thể rút trích từ tên lớp, nên bổ sung thêm mô tả chi tiết từ LLM không mang lại lợi ích đáng kể cho giai đoạn phát hiện.

    \item \textbf{YOLO-World đơn lẻ có MAE cao hơn baseline} (19,03 so với 17,99). Nguyên nhân chính là mô hình nhẹ này, dù nhanh hơn đáng kể, có năng lực nhận dạng ngữ nghĩa yếu hơn GroundingDINO, dẫn đến nhiễu phát hiện (detection noise) cao hơn ở giai đoạn trích xuất mẫu. Cụ thể, YOLO-World sử dụng chiến lược Prompt-then-Detect (mã hóa ngoại tuyến text prompt), nên không có cơ chế tương tác thị giác-ngôn ngữ (vision-language interaction) tại thời điểm suy luận, dẫn đến khả năng phân biệt đối tượng mục tiêu và nhiễu nền kém hơn so với GroundingDINO.

    \item \textbf{YOLO-World + Rich Prompt đạt cân bằng tốt nhất:} Rich Prompt bù đắp được hạn chế nhận dạng của YOLO-World, giảm MAE từ 19,03 xuống 17,91 ($\Delta = 1{,}12$), gần ngang VA-Count gốc (17,99) nhưng với chi phí tính toán thấp hơn đáng kể (xem Bảng~\ref{tab:demo_time}). Mức cải thiện $\Delta = 1{,}12$ trên YOLO-World so với $\Delta = 0{,}19$ trên GroundingDINO khẳng định rằng: \emph{bộ phát hiện càng yếu về ngữ nghĩa, Rich Prompt càng mang lại lợi ích lớn}. Điều này gợi mở một chiến lược thiết kế hệ thống: kết hợp bộ phát hiện nhẹ với prompting thông minh có thể đạt hiệu quả tương đương bộ phát hiện nặng, nhưng với tốc độ nhanh hơn nhiều lần.
\end{itemize}

\subsection{Đặt kết quả trong bức tranh toàn cảnh}
\label{subsec:sota-comparison}

Để đánh giá đúng vị trí của hệ thống, Bảng~\ref{tab:sota} đặt kết quả bốn cấu hình của nhóm cạnh các phương pháp tiêu biểu trên cùng FSC-147 test split. Ba nhóm được phân theo loại đầu vào mẫu: \textit{Reference-less} (không cần mẫu), \textit{Few-shot} (mẫu hình ảnh từ người dùng), và \textit{Zero-shot / Text} (chỉ dùng tên lớp văn bản - cùng điều kiện với hệ thống của nhóm).

\begin{table}[H]
\centering
\small
\renewcommand{\arraystretch}{1.2}
\begin{adjustbox}{max width=\textwidth}
\begin{tabular}{|l|l|l|c|c|}
\hline
\rowcolor[HTML]{EFEFEF}
\textbf{Nhóm} & \textbf{Phương pháp} & \textbf{Venue} & \textbf{MAE}$\downarrow$ & \textbf{RMSE}$\downarrow$ \\
\hline
\textit{Ref-less}  & FamNet~\cite{famnet}               & CVPR'21  & 32{,}27          & 131{,}46 \\
                   & RCC$^\dagger$                      & CVPR'22  & 17{,}12          & 104{,}53 \\
                   & CounTR~\cite{countr}               & BMVC'22  & \textbf{14{,}71} & 106{,}87 \\
                   & LOCA$^\dagger$                     & ICCV'23  & 16{,}22          & 103{,}96 \\
\hline
\textit{Few-shot}  & FamNet~\cite{famnet}               & CVPR'21  & 22{,}56          & 101{,}54 \\
                   & CounTR~\cite{countr}               & BMVC'22  & 11{,}95          & 91{,}23  \\
                   & LOCA$^\dagger$                     & ICCV'23  & \textbf{10{,}97} & \textbf{56{,}97} \\
\hline
\textit{Zero-shot} & ZSC~\cite{zeroshot_counting_xu}    & CVPR'23  & 22{,}09          & 115{,}17 \\
                   & CLIP-Count$^\dagger$               & MM'23    & 17{,}19          & 103{,}44 \\
                   & VLCount$^\dagger$                  & AAAI'24  & 17{,}05          & 106{,}16 \\
                   & VA-Count~\cite{vacount} (gốc)      & ECCV'24  & 17{,}88          & 129{,}31 \\
\rowcolor[HTML]{E8F5E9}
                   & VA-Count (nhóm)$^\ddagger$         & {-}    & 17{,}99          & 129{,}39 \\
\rowcolor[HTML]{E8F5E9}
                   & \quad + Rich Prompt$^\ddagger$     & {-}    & \textbf{17{,}80} & 129{,}69 \\
\rowcolor[HTML]{E8F5E9}
                   & \quad + YOLO-World$^\ddagger$      & {-}    & 19{,}03          & 131{,}55 \\
\rowcolor[HTML]{E8F5E9}
                   & \quad + YOLO-World + RP$^\ddagger$ & {-}    & 17{,}91          & 130{,}98 \\
\hline
\end{tabular}
\end{adjustbox}
\caption{So sánh với các phương pháp tiêu biểu trên FSC-147 \textbf{test split}. In đậm: tốt nhất trong mỗi nhóm. Nền xanh lá~($^\ddagger$): bốn cấu hình thực nghiệm của nhóm. Dấu~$^\dagger$: số liệu trích từ bảng tổng hợp trong~\cite{richprompts}, không có bibitem riêng trong tài liệu này.}
\label{tab:sota}
\end{table}

Một số nhận xét từ bảng so sánh:

\begin{itemize}
    \item \textbf{Khoảng cách giữa zero-shot và few-shot.} Phương pháp few-shot tốt nhất (LOCA, MAE $= 10{,}97$) vẫn hơn VA-Count baseline của nhóm ($17{,}99$) khoảng $7$ điểm MAE. Khoảng cách này là cái giá phải trả cho việc loại bỏ sự can thiệp thủ công: khi người dùng cung cấp ảnh mẫu chất lượng cao, bộ đếm hoạt động tốt hơn đáng kể. Đây là giới hạn cơ bản của hướng zero-shot, không phải riêng của VA-Count.
    \item \textbf{Hiện thực trung thực.} Nhóm tái hiện VA-Count đạt MAE $= 17{,}99$, chỉ chênh $+0{,}11$ so với kết quả gốc ($17{,}88$) từ~\cite{vacount}. Sự nhất quán này xác nhận rằng các kết quả thực nghiệm của nhóm đáng tin cậy và có thể so sánh trực tiếp với tài liệu gốc.
    \item \textbf{Bộ đếm là nhân tố quyết định.} CounTR ở chế độ reference-less đạt MAE $= 14{,}71$, tốt hơn mọi phương pháp zero-shot text-based trong bảng. Nhất quán với phân tích tại Mục~\ref{sec:results}: khoảng dao động MAE giữa bốn cấu hình của nhóm chỉ $1{,}23$ điểm cho thấy \emph{chất lượng kiến trúc bộ đếm là nhân tố quyết định}, không phải chiến lược chọn mẫu.
\end{itemize}

\subsection{Tốc độ suy luận khi demo}

\begin{table}[H]
\centering
\renewcommand{\arraystretch}{1.2}
\begin{tabular}{|l|c|}
\hline
\rowcolor[HTML]{EFEFEF}
\textbf{Phương pháp} & \textbf{Thời gian trung bình (s/ảnh)} \\
\hline
VA-Count & 1,4710 \\
\hline
VA-Count + Rich Prompt & 5,7578 \\
\hline
VA-Count với YOLO-World & \textbf{0,6006} \\
\hline
VA-Count với YOLO-World + Rich Prompt & 2,4054 \\
\hline
\end{tabular}
\caption{Thời gian suy luận (inference) trung bình trên 1 ảnh trong môi trường demo (RTX~4060).}
\label{tab:demo_time}
\end{table}

Trên cấu hình demo, YOLO-World giảm thời gian xuống $0{,}60$ s/ảnh, nhanh gấp 2,5 lần VA-Count gốc. Khi bật Rich Prompt, hệ thống giữ ở mức $2{,}4$ s/ảnh, vẫn chấp nhận được với tương tác người dùng. Phương án \textbf{YOLO-World + Rich Prompt} đạt cân bằng tốt nhất giữa độ chính xác (MAE $17{,}91$) và độ trễ ($2{,}4$ s/ảnh), và đây là phương án mặc định mà ứng dụng demo của nhóm sử dụng (Chương~\ref{ch:demo}).

\paragraph{Nhận xét về đánh đổi tốc độ-chính xác.}
Kết hợp Bảng~\ref{tab:counting_perf} và Bảng~\ref{tab:demo_time}, có thể thấy rõ sự đánh đổi giữa tốc độ và chất lượng đếm. VA-Count gốc với GroundingDINO đạt MAE tốt nhất (17,99) nhưng bị hạn chế bởi độ trễ ($1{,}47$ s/ảnh), đặc biệt khi bật Rich Prompt thì lên tới $5{,}76$ s/ảnh, gần như không khả thi cho ứng dụng thời gian thực. YOLO-World cải thiện tốc độ lên mức $0{,}60$ s/ảnh (nhanh gấp 2,5 lần), đổi lại MAE tăng lên 19,03. Tuy nhiên, khi kết hợp với Rich Prompt, YOLO-World đạt MAE $17{,}91$ (chỉ thua baseline $0{,}08$) với thời gian $2{,}4$ s/ảnh. Kết quả này chứng tỏ phương án đề xuất đạt được sự cân bằng tốt giữa hiệu năng tính toán và chất lượng dự đoán, phù hợp cho triển khai thực tế.
Hình~\ref{fig:latency_accuracy} trực quan hóa sự đánh đổi này: mỗi điểm trên biểu đồ ứng với một cấu hình, trục hoành là độ trễ (s/ảnh), trục tung là MAE. Cấu hình lý tưởng nằm ở góc dưới-trái (nhanh \emph{và} chính xác). YOLO-World + Rich Prompt đạt vị trí cân bằng tốt nhất, gần nhóm chính xác cao nhưng ở vùng tốc độ chấp nhận được.

\begin{figure}[H]
\centering
\begin{tikzpicture}
\begin{axis}[
    width=0.85\textwidth,
    height=0.55\textwidth,
    xlabel={Độ trễ suy luận (s/ảnh)},
    ylabel={MAE $\downarrow$},
    xmin=0, xmax=6.5,
    ymin=17.0, ymax=20.0,
    grid=major,
    grid style={dashed,gray!40},
    legend pos=north east,
    legend style={font=\small, draw=gray!50},
    nodes near coords,
    every node near coord/.append style={font=\scriptsize, anchor=south west, xshift=2pt},
    point meta=explicit symbolic,
    scatter/classes={
        a={mark=*,draw=blue!80,fill=blue!50,mark size=3.5pt},
        b={mark=square*,draw=red!80,fill=red!50,mark size=3.5pt},
        c={mark=triangle*,draw=teal!80,fill=teal!50,mark size=3.5pt},
        d={mark=diamond*,draw=orange!80,fill=orange!60,mark size=4pt}
    },
]
\addplot[scatter, only marks, scatter src=explicit symbolic]
    table[x=latency, y=mae, meta=class] {
        latency  mae    class  label
        0.6006   19.03  a      {YOLO-World}
        1.4710   17.99  b      {VA-Count}
        2.4054   17.91  d      {YOLO+RP}
        5.7578   17.80  c      {GDino+RP}
    };
\legend{YOLO-World, VA-Count, YOLO-World + RP, GDino + RP}
% Arrow highlighting Pareto-optimal region
\draw[->,thick,dashed,gray!60] (axis cs:3.0,19.5) -- (axis cs:0.8,17.5)
    node[pos=0.4,above,sloped,font=\scriptsize\itshape,text=gray]{tối ưu};
\end{axis}
\end{tikzpicture}
\caption{Đánh đổi tốc độ--chính xác (Latency vs.\ Accuracy) trên FSC-147 test split. Mũi tên nét đứt chỉ hướng tối ưu (nhanh hơn và chính xác hơn). YOLO-World + Rich Prompt (hình thoi cam) đạt vị trí cân bằng tốt nhất giữa hai tiêu chí.}
\label{fig:latency_accuracy}
\end{figure}

\subsection{Phân tích MAE theo mật độ đối tượng}
\label{subsec:mae-vs-density}

Để hiểu sâu hơn ảnh hưởng của mật độ đối tượng (object density) đến chất lượng đếm, nhóm chia tập test FSC-147 (1\,190 ảnh) thành các nhóm theo mật độ (số đối tượng trên mỗi megapixel) và tính MAE trung bình cho mỗi nhóm. Hình~\ref{fig:mae_density} cho thấy MAE tăng rõ rệt khi mật độ tăng: ở vùng mật độ thấp ($<50$ obj/Mpx), cả bốn cấu hình đều đạt MAE dưới 10; tuy nhiên ở vùng mật độ cao ($>500$ obj/Mpx), MAE vượt 50, phản ánh đúng hiện tượng \emph{đếm thiếu ở mật độ dày đặc} đã phân tích tại Mục~\ref{sec:error-analysis}. Đáng chú ý, YOLO-World không có Rich Prompt (đường xanh lam) luôn có MAE cao hơn các cấu hình khác ở mọi nhóm mật độ, trong khi YOLO-World + Rich Prompt (cam) bám sát VA-Count baseline, xác nhận vai trò bù đắp ngữ nghĩa của Rich Prompt.

\begin{figure}[H]
\centering
\begin{tikzpicture}
\begin{axis}[
    width=0.88\textwidth,
    height=0.55\textwidth,
    xlabel={Mật độ đối tượng (obj/Mpx)},
    ylabel={MAE trung bình $\downarrow$},
    xmode=log,
    log basis x=10,
    xmin=5, xmax=5000,
    ymin=0,
    grid=major,
    grid style={dashed,gray!40},
    legend pos=north west,
    legend style={font=\small, draw=gray!50},
    legend cell align={left},
    mark options={solid},
]
\addplot[blue!80, thick, mark=*,     mark size=2.5pt] table[x=bin_mid, y=baseline]   {figures/mae_vs_density.dat};
\addplot[red!80,  thick, mark=square*,mark size=2.5pt] table[x=bin_mid, y=richprompt] {figures/mae_vs_density.dat};
\addplot[teal!80, thick, mark=triangle*,mark size=2.5pt] table[x=bin_mid, y=yolo]     {figures/mae_vs_density.dat};
\addplot[orange!80,thick,mark=diamond*,mark size=3pt]  table[x=bin_mid, y=yolo_rp]   {figures/mae_vs_density.dat};
\legend{VA-Count, +Rich Prompt, +YOLO-World, +YOLO+RP}
\end{axis}
\end{tikzpicture}
\caption{MAE trung bình theo mật độ đối tượng (logarithmic scale) trên FSC-147 test split. Mỗi điểm đại diện cho MAE trung bình của tất cả ảnh thuộc một khoảng mật độ. Mật độ cao đi kèm sai số lớn hơn đáng kể ở mọi cấu hình.}
\label{fig:mae_density}
\end{figure}

\subsection{Phân tích MAE theo kích thước đối tượng}
\label{subsec:mae-vs-size}

Tương tự, nhóm phân tích ảnh hưởng của kích thước đối tượng (đo bằng diện tích trung bình của hộp giới hạn mẫu, tính theo pixel$^2$) đến MAE. Hình~\ref{fig:mae_size} cho thấy xu hướng \emph{ngược lại} so với mật độ: đối tượng nhỏ ($<500$ px$^2$) có MAE cao nhất vì chúng thường xuất hiện với mật độ dày đặc và dễ bị bỏ sót bởi bản đồ mật độ; đối tượng lớn ($>6\,000$ px$^2$) cho MAE thấp nhất. Khoảng cách giữa YOLO-World (không có Rich Prompt) và các cấu hình khác ở nhóm đối tượng nhỏ là lớn nhất ($\sim$4--5 điểm MAE), cho thấy YOLO-World gặp khó khăn đặc biệt với các đối tượng nhỏ do thiếu cơ chế tương tác thị giác-ngôn ngữ chi tiết. Rich Prompt giúp thu hẹp khoảng cách này đáng kể.

\begin{figure}[H]
\centering
\begin{tikzpicture}
\begin{axis}[
    width=0.88\textwidth,
    height=0.55\textwidth,
    xlabel={Kích thước đối tượng trung bình (px$^2$)},
    ylabel={MAE trung bình $\downarrow$},
    xmode=log,
    log basis x=10,
    xmin=100, xmax=30000,
    ymin=0,
    grid=major,
    grid style={dashed,gray!40},
    legend pos=north east,
    legend style={font=\small, draw=gray!50},
    legend cell align={left},
    mark options={solid},
]
\addplot[blue!80, thick, mark=*,     mark size=2.5pt] table[x=bin_mid, y=baseline]   {figures/mae_vs_size.dat};
\addplot[red!80,  thick, mark=square*,mark size=2.5pt] table[x=bin_mid, y=richprompt] {figures/mae_vs_size.dat};
\addplot[teal!80, thick, mark=triangle*,mark size=2.5pt] table[x=bin_mid, y=yolo]     {figures/mae_vs_size.dat};
\addplot[orange!80,thick,mark=diamond*,mark size=3pt]  table[x=bin_mid, y=yolo_rp]   {figures/mae_vs_size.dat};
\legend{VA-Count, +Rich Prompt, +YOLO-World, +YOLO+RP}
\end{axis}
\end{tikzpicture}
\caption{MAE trung bình theo kích thước đối tượng (logarithmic scale) trên FSC-147 test split. Đối tượng nhỏ (diện tích $<500$ px$^2$) có MAE cao nhất; kích thước tăng thì MAE giảm. YOLO-World không có Rich Prompt chịu sai số lớn nhất trên đối tượng nhỏ.}
\label{fig:mae_size}
\end{figure}

\subsection{So sánh định tính giữa VA-Count và YOLO-World}
\label{subsec:qualitative-comparison}

Để minh họa rõ hơn sự khác biệt giữa hai bộ phát hiện, nhóm so sánh trực quan kết quả trên cùng một ảnh xuất hiện trong cả hai thí nghiệm.

\begin{figure}[H]
    \centering
    \begin{subfigure}{0.48\textwidth}
        \centering
        \figureorplaceholder{figures/exp2_result_711.jpg}{\linewidth}{VA-Count (GroundingDINO) trên ảnh 711.}
        \caption{VA-Count (GroundingDINO)}
    \end{subfigure}\hfill
    \begin{subfigure}{0.48\textwidth}
        \centering
        \figureorplaceholder{figures/exp3_best_711.png}{\linewidth}{YOLO-World trên ảnh 711 (best case).}
        \caption{YOLO-World (best case)}
    \end{subfigure}
    \caption{So sánh kết quả trên ảnh \texttt{711}: YOLO-World chọn được mẫu (exemplar) chính xác, bản đồ mật độ tập trung tốt vào đối tượng mục tiêu, cho kết quả đếm tương đương hoặc tốt hơn VA-Count gốc trên trường hợp này. Ảnh trái: \texttt{experiments/exp2/visualizations/result\_711.jpg}; ảnh phải: \texttt{experiments/exp3/visualizations\_test/best\_2\_711.png}.}
    \label{fig:compare_711}
\end{figure}

\begin{figure}[H]
    \centering
    \begin{subfigure}{0.48\textwidth}
        \centering
        \figureorplaceholder{figures/exp2_result_716.jpg}{\linewidth}{VA-Count (GroundingDINO) trên ảnh 716.}
        \caption{VA-Count (GroundingDINO)}
    \end{subfigure}\hfill
    \begin{subfigure}{0.48\textwidth}
        \centering
        \figureorplaceholder{figures/exp3_worst_716.png}{\linewidth}{YOLO-World trên ảnh 716 (worst case).}
        \caption{YOLO-World (worst case)}
    \end{subfigure}
    \caption{So sánh kết quả trên ảnh \texttt{716}: trường hợp YOLO-World gặp khó khăn. Do năng lực nhận dạng ngữ nghĩa yếu hơn, YOLO-World chọn mẫu kém chính xác hơn, dẫn đến bản đồ mật độ phân tán và sai số đếm lớn hơn so với VA-Count gốc. Ảnh trái: \texttt{experiments/exp2/visualizations/result\_716.jpg}; ảnh phải: \texttt{experiments/exp3/visualizations\_test/worst\_3\_716.png}.}
    \label{fig:compare_716}
\end{figure}

Hình~\ref{fig:compare_711} cho thấy trên ảnh \texttt{711}, YOLO-World chọn mẫu chính xác và cho kết quả đếm tốt, chứng minh rằng bộ phát hiện nhẹ vẫn hoàn toàn có thể cạnh tranh với GroundingDINO khi đối tượng có đặc trưng thị giác rõ ràng và ít bị nhầm lẫn với nền. Ngược lại, Hình~\ref{fig:compare_716} minh họa trường hợp YOLO-World gặp khó khăn: khi đối tượng có hình dạng hoặc kết cấu (texture) tương tự vật thể nền, việc thiếu cơ chế tương tác thị giác-ngôn ngữ tại thời điểm suy luận khiến YOLO-World chọn sai mẫu, dẫn đến bản đồ mật độ bị phân tán và sai số đếm tăng đáng kể. Đây chính là lý do Rich Prompt mang lại cải thiện lớn ($\Delta$MAE $= 1{,}12$) cho YOLO-World: mô tả ngữ nghĩa chi tiết từ LLM giúp bù đắp khả năng phân biệt mà YOLO-World thiếu.

\subsection{Minh họa kết quả đếm đúng}
\label{subsec:vis-correct}

Ngoài các trường hợp so sánh ở trên, hệ thống còn cho thấy khả năng đếm chính xác cao trên nhiều loại đối tượng đa dạng. Hình~\ref{fig:vis1},~\ref{fig:vis2} và~\ref{fig:vis3} minh họa ba ví dụ tiêu biểu, trong đó bên trái là ảnh gốc kèm positive exemplars (xanh lá) và negative exemplars (đỏ), bên phải là density map với overlay dự đoán.

\begin{figure}[H]
    \centering
    \figureorplaceholder{figures/vis1.png}
    {0.95\textwidth}{Kết quả đếm đúng trên ảnh bi thủy tinh: GT=11, Predicted=11.2.}
    \caption{Kết quả đếm trên ảnh bi thủy tinh: đối tượng có bề mặt phản chiếu gây khó khăn cho trích xuất đặc trưng, nhưng mô hình vẫn cho dự đoán chính xác. GT $= 11$, dự đoán $= 11{,}2$ (sai số $1{,}8\%$).}
    \label{fig:vis1}
\end{figure}

\begin{figure}[H]
    \centering
    \figureorplaceholder{figures/vis2.png}{0.95\textwidth}{Kết quả đếm đúng trên ảnh cừu: GT=32, Predicted=32.2.}
    \caption{Kết quả đếm trên ảnh cừu chụp từ trên cao: dù đối tượng có bóng đổ và sắp xếp không đều, hệ thống vẫn định vị chính xác từng con cừu trên bản đồ mật độ. GT $= 32$, dự đoán $= 32{,}2$ (sai số $0{,}6\%$).}
    \label{fig:vis2}
\end{figure}

\begin{figure}[H]
    \centering
    \figureorplaceholder{figures/vis3.png}
    {0.95\textwidth}{Kết quả đếm đúng trên ảnh táo: GT=33, Predicted=32.67.}
    \caption{Kết quả đếm trên ảnh chứa táo: hệ thống chọn được positive exemplars chính xác (xanh lá), negative exemplar (đỏ) loại bỏ nền. Bản đồ mật độ tập trung đúng vị trí từng quả. GT $= 33$, dự đoán $= 32{,}67$ (sai số $1\%$).}
    \label{fig:vis3}
\end{figure}

\section{Phân tích lỗi sai}
\label{sec:error-analysis}
Mặc dù các giải pháp đề xuất đã đạt được hiệu suất tổng thể khá tốt trên tập dữ liệu FSC147, hệ thống vẫn bộc lộ những hạn chế nhất định khi đối mặt với một vài tình huống. Việc nhận diện thẳng thắn các điểm yếu này là cơ sở quan trọng để hiểu rõ giới hạn vận hành của mô hình. Qua thực nghiệm, nhóm đã ghi nhận sai số thường tập trung vào hai nhóm nguyên nhân chính sau đây: \textbf{Mật độ dày đặc} và \textbf{Sự phân mảnh đối tượng.}

\subsection{Mật độ dày đặc}

Trong các trường hợp đối tượng được sắp xếp chồng
lấn nghiêm trọng hoặc kích thước đối tượng quá nhỏ và sát nhau (ví dụ hộp bút màu hoặc hạt cườm), ranh giới thị giác (visual boundary) giữa các vật thể trở nên mờ nhạt, Sự hòa trộn đặc trưng này khiến mô hình gặp khó khăn trong việc tách biệt đâu là vật thể, đâu là nền, từ đó làm bản đồ mật độ suy giảm giá trị và gây ra lỗi \textbf{đếm thiếu} đáng kể.

\begin{figure}[H]
    \centering
    \figureorplaceholder{figures/failure_dense.png}{0.95\textwidth}{Trường hợp đếm thiếu khi mật độ dày đặc.}
    \caption{Minh họa sai số đếm thiếu trên ảnh có mật độ rất dày: GT $= 548$, dự đoán $\approx 283{,}55$, sai số tuyệt đối $264{,}45$.}
    \label{fig:fail_dense}
\end{figure}
Hình~\ref{fig:fail_dense_pens} minh họa thêm một trường hợp điển hình: hộp bút màu với $371$ đối tượng xếp chồng lấn nghiêm trọng. Mô hình chỉ dự đoán được $\approx 251{,}55$, cho thấy sự suy giảm đáng kể của bản đồ mật độ khi ranh giới giữa các đối tượng gần như biến mất.

\begin{figure}[H]
    \centering
    \figureorplaceholder{figures/dem-sai-do-mat-do-day.png}{0.95\textwidth}{Trường hợp đếm sai do mật độ dày: hộp bút màu.}
    \caption{Minh họa cho trường hợp sai số nhiều do mật độ dày: ảnh gốc (trái) với positive exemplars (xanh lá) và negative exemplar (đỏ), density map (phải) cho thấy nhiều vùng bị hòa trộn. GT $= 371$, dự đoán $\approx 251{,}55$, sai số tuyệt đối $119{,}45$.}
    \label{fig:fail_dense_pens}
\end{figure}

\subsection{Phân mảnh đối tượng}

Các vật thể có cấu trúc chia tách (cửa sổ nhiều ô, kính mắt) hoặc kết cấu bề mặt phức tạp bị nhận diện nhầm thành nhiều thực thể độc lập (như từng ô kính, tròng kính) thay vì một khối thống nhất, tạo nhiều điểm cực đại cục bộ (local maxima) trên cùng một đối tượng và gây sai số \textbf{đếm thừa}.


\begin{figure}[H]
    \centering
    \figureorplaceholder{figures/failure_fragment.png}{0.85\textwidth}{Trường hợp đếm thừa do phân mảnh đối tượng (cửa sổ nhiều ô).}
    \caption{Failure case: phân mảnh đối tượng. Dãy cửa sổ nhiều ô bị tách thành nhiều cực đại cục bộ trên bản đồ mật độ (GT $= 13$, dự đoán $\approx 27{,}13$). Ví dụ định tính trích từ phân tích nội bộ của nhóm trên FSC-147.}
    \label{fig:fail_fragment}
\end{figure}

\subsection{Sai số mang tính phương pháp}

Mô hình dùng một bộ trích xuất đặc trưng (feature extractor) kiểu Masked Autoencoder + bản đồ mật độ. Quan sát thực tế cho thấy mô hình \emph{vẫn đưa ra ước lượng} ngay cả khi không lấy được mẫu hợp lệ, tuy nhiên sai số rất lớn (ví dụ ground truth $371$, dự đoán $251{,}55$). Hiện tượng này gợi ý rằng mô hình có xu hướng dựa vào các đặc trưng toàn cục (global feature) (mật độ, kết cấu, bố cục không gian) thay vì thực sự condition theo mẫu. Đây là một dạng \textbf{thiên lệch dữ liệu huấn luyện (training data bias)} khiến vai trò của mẫu trong quá trình suy luận bị suy giảm.\\[1em]
Nguyên nhân có thể xuất phát từ việc mô hình học được các đặc trưng tương quan với mật độ đối tượng chung (như kết cấu lặp lại, bố cục đồng đều) trong quá trình huấn luyện trên FSC-147, một bộ dữ liệu có phân bố mật độ khá tập trung. Đây là hạn chế cơ bản của phương pháp đếm dựa trên bản đồ mật độ, đặc biệt trong các kịch bản mật độ cao hoặc phân bố phức tạp, và cần được giải quyết trong các phiên bản kế tiếp (xem Chương~\ref{ch:conclusion}).

%==================================================================
%                          CHƯƠNG 4
%==================================================================
\chapter{DEMO VÀ ĐÁNH GIÁ HỆ THỐNG}
\label{ch:demo}

\begin{flushleft}
    \textbf{Mã nguồn dự án:} \href{https://github.com/paht2005/CS338.Q21_Zero-shot-Object-Coutning-with-Good-Examplers}{GitHub Repository} \\
    \textbf{Video Demo:} cập nhật trong file \texttt{README.md} (mục \emph{Demo \& Outputs}).
\end{flushleft}
Nhóm cung cấp một ứng dụng demo dạng Streamlit~\cite{streamlit} (\texttt{code/source-code/demo\_app\_advanced.py}) cùng các script hỗ trợ:

\begin{itemize}
    \item \texttt{demo\_app\_advanced.py}: ứng dụng Streamlit nâng cao, cho phép tùy chọn bộ trích xuất (GroundingDINO / YOLO-World), bật/tắt Rich Prompt, hiển thị bản đồ mật độ, hộp giới hạn dương/âm và so sánh với ground-truth (mặc định: \textbf{YOLO-World + Rich Prompt}).
    \item \texttt{demo\_inference.py}: thư viện suy luận dùng chung (load model, expand prompt, pipeline tối giản); được \texttt{demo\_app\_advanced.py} import.
    \item \texttt{demo\_pipeline\_advanced.py}, \texttt{demo\_visualization.py}: module pipeline đầy đủ và các hàm vẽ overlay được dùng bên trong ứng dụng Streamlit.
\end{itemize}

\section{Cài đặt và chạy demo}

\begin{lstlisting}[style=bashstyle]
cd code/source-code
python -m venv .venv && source .venv/bin/activate
pip install -r requirements.txt

# 1. Set Gemini API key (DO NOT commit)
cp .env.example .env
# edit .env and fill in GEMINI_API_KEY

# 2. Run the Streamlit demo
streamlit run demo_app_advanced.py
\end{lstlisting}

\section{Giao diện và tương tác}
Người dùng tải lên một ảnh tự nhiên, nhập tên lớp tiếng Anh (ví dụ \texttt{oranges}), và nhận về:
\begin{itemize}
    \item Tổng số đối tượng đếm được.
    \item Bản đồ mật độ overlay trên ảnh gốc.
    \item Các hộp giới hạn (bbox) dương/âm đã được EEM chọn.
\end{itemize}

\begin{figure}[H]
    \centering
    \figureorplaceholder{figures/demo_streamlit.png}{0.95\textwidth}{Đầu ra ba ô (ảnh gốc với positive exemplars, predicted density map, overlay) cùng bảng so sánh số lượng dự đoán và ground truth.}
    \caption{Đầu ra của hệ thống cho ảnh \texttt{346.jpg} (lớp \emph{strawberries}): ảnh gốc kèm 3 mẫu dương (trái), bản đồ mật độ dự đoán (giữa), overlay heatmap lên ảnh gốc (phải) và bảng tổng hợp số đếm. Ground truth $= 11$, mô hình dự đoán $\approx 11{,}26$ (sai số tuyệt đối $0{,}26$). Visualization từ \texttt{experiments/exp3/visualizations\_test/best\_1\_346.png}; ứng dụng demo Streamlit dùng cùng pipeline để render.}
    \label{fig:demo}
\end{figure}

\section{Đánh giá định tính}
Kiểm thử trên các ảnh khó từ FSC-147 và một số ảnh ngoài tập (chợ trái cây, ảnh chăn nuôi) cho thấy:
\begin{itemize}
    \item Cấu hình \textbf{YOLO-World + Rich Prompt} cho phản hồi nhanh ($\sim 2$ s) và chấp nhận được trên CPU/GPU laptop.
    \item Khi tên lớp rõ ràng và đối tượng phân tách tốt, kết quả gần với ground truth.
    \item Trên các ảnh dày đặc hoặc tên lớp mơ hồ, hệ thống bộc lộ đúng các lỗi đã phân loại ở Mục~\ref{sec:error-analysis}.
\end{itemize}

%==================================================================
%                          CHƯƠNG 5
%==================================================================
\chapter{KẾT LUẬN VÀ HƯỚNG PHÁT TRIỂN}
\label{ch:conclusion}

\section{Kết luận}

Đồ án này hiện thực lại và mở rộng hệ thống VA-Count~\cite{vacount} theo hai hướng độc lập, đồng thời đánh giá toàn diện bốn cấu hình trên FSC-147. Kết quả tái hiện (MAE $= 17{,}99$) chỉ chênh $+0{,}11$ so với bài báo gốc ($17{,}88$), xác nhận tính trung thực của hiện thực. Trong bức tranh toàn cảnh (Bảng~\ref{tab:sota}), cả bốn cấu hình đều nằm trong nhóm dẫn đầu của zero-shot text-based methods, đồng thời cho thấy khoảng cách $\approx 7$ điểm MAE so với few-shot tốt nhất là giới hạn tự nhiên khi thiếu mẫu chất lượng cao từ người dùng.\\[1em]
Phát hiện quan trọng nhất của đồ án là hiện tượng mà nhóm đặt tên \textbf{Giả thuyết bão hòa ngữ nghĩa} (\textit{Semantic Saturation Hypothesis}): bộ phát hiện có kiến trúc hợp nhất thị giác--ngôn ngữ mạnh (GroundingDINO) đã khai thác gần hết lượng thông tin ngữ nghĩa có thể rút trích từ text đầu vào, nên bổ sung mô tả chi tiết từ LLM chỉ giảm thêm $\Delta\mathrm{MAE} = 0{,}19$. Ngược lại, bộ phát hiện nhẹ hơn (YOLO-World) hưởng lợi đáng kể ($\Delta\mathrm{MAE} = 1{,}12$), đến mức YOLO-World + Rich Prompt ($17{,}91$) gần như bắt kịp VA-Count baseline ($17{,}99$) dù YOLO-World đơn lẻ chỉ đạt $19{,}03$. Giả thuyết này có hàm ý thiết kế rõ ràng: \emph{đầu tư vào prompt engineering tạo ra lợi nhuận cận biên giảm dần khi bộ phát hiện ngày càng mạnh về ngữ nghĩa}; hướng cải tiến tiềm năng hơn là nâng cấp kiến trúc bộ đếm (Counter/NSM) - thành phần chung của cả bốn cấu hình.

\section{Hướng phát triển}

Trên cơ sở các phân tích lỗi sai (Mục~\ref{sec:error-analysis}) và các lựa chọn cài đặt nêu trên, nhóm đề xuất các hướng phát triển sau:

\begin{itemize}
    \item \textbf{Khử phân mảnh đối tượng.} Áp dụng các kỹ thuật hậu xử lý (post-processing) (graph-based clustering, watershed) hoặc thêm head phụ dự đoán \textit{bản đồ mật độ nhận biết thực thể (instance-aware density)} để giảm đếm thừa cho các vật thể có cấu trúc chia tách.
    \item \textbf{Mật độ dày đặc.} Bổ sung kiến trúc đa độ phân giải (multi-resolution) hoặc suy luận theo ô (tile-based inference) để giữ chi tiết tại biên các vật thể sát nhau.
    \item \textbf{Khử thiên lệch đếm không cần mẫu.} Thử nghiệm các phương pháp đếm dựa trên đối sánh trực tiếp (matching) (CountR~\cite{countr}) thay cho bản đồ mật độ kiểu MAE, hoặc thêm regularizer ép mô hình bằng 0 khi mẫu không hợp lệ.
    \item \textbf{Giảm chi phí tính toán của Rich Prompt.} Thời gian xử lý $\approx 25$ giờ cho cấu hình GroundingDINO + Rich Prompt xuất phát chủ yếu từ việc gọi Gemini API \emph{riêng lẻ cho từng ảnh} trong bước sinh negative prompt (phụ thuộc nội dung ảnh cụ thể). Tuy nhiên, \emph{positive prompt} chỉ phụ thuộc vào tên lớp, không phụ thuộc vào ảnh cụ thể. Do đó, với FSC-147 có 147 lớp, toàn bộ positive prompts có thể được \textit{precompute} một lần ($\approx 147$ API calls, vài phút) và tái sử dụng cho mọi ảnh trong cùng lớp. Negative prompt có thể được xấp xỉ bằng class-level description (mô tả nền chung cho mỗi lớp) thay vì per-image call. Kết hợp hai kỹ thuật này có thể đưa tổng thời gian chuẩn bị dữ liệu của cấu hình GroundingDINO + Rich Prompt về mức tương đương YOLO-World + Rich Prompt ($\approx 4$ giờ), mở ra khả năng thực nghiệm nhanh hơn nhiều lần.
    \item \textbf{Triển khai sản phẩm.} Đóng gói pipeline trong Docker và triển khai trên dịch vụ HuggingFace Spaces / Vercel để tăng khả năng tiếp cận.
\end{itemize}

%==================================================================
%                          BIBLIOGRAPHY
%==================================================================

\addcontentsline{toc}{chapter}{TÀI LIỆU THAM KHẢO}

\begin{thebibliography}{99}

\bibitem{vacount}
H.~Zhu, J.~Yuan, Z.~Yang, Y.~Guo, Z.~Wang, X.~Zhong, and S.~He,
``Zero-shot object counting with good exemplars,''
in \emph{European Conference on Computer Vision (ECCV)}, 2024.
[Online]. Available:
\url{https://www.ecva.net/papers/eccv_2024/papers_ECCV/papers/00812.pdf}

\bibitem{richprompts}
H.~Zhu, S.~Li, J.~Yuan, Z.~Yang, Y.~Guo, W.~Liu, X.~Zhong, and S.~He,
``Expanding zero-shot object counting with rich prompts,''
arXiv preprint arXiv:2505.15398, 2025.
[Online]. Available:
\url{https://arxiv.org/abs/2505.15398}

\bibitem{countr}
C.~Liu, Y.~Zhong, A.~Zisserman, and W.~Xie,
``CounTR: Transformer-based generalised visual counting,''
in \emph{British Machine Vision Conference (BMVC)}, 2022.
[Online]. Available:
\url{https://pdf-files.bmvc2022.org/0370.pdf}

\bibitem{yoloworld}
T.~Cheng, L.~Song, Y.~Ge, W.~Liu, X.~Wang, and Y.~Shan,
``YOLO-World: Real-time open-vocabulary object detection,''
in \emph{IEEE/CVF Conference on Computer Vision and Pattern
Recognition (CVPR)}, 2024.
[Online]. Available:
\url{https://openaccess.thecvf.com/content/CVPR2024/papers/Cheng_YOLO-World_Real-Time_Open-Vocabulary_Object_Detection_CVPR_2024_paper.pdf}

\bibitem{groundingdino}
S.~Liu \emph{et~al.},
``Grounding DINO: Marrying DINO with grounded pre-training
for open-set object detection,''
in \emph{European Conference on Computer Vision (ECCV)}, 2024.
[Online]. Available:
\url{https://arxiv.org/abs/2303.05499}

\bibitem{clip}
A.~Radford \emph{et~al.},
``Learning transferable visual models from natural language
supervision (CLIP),''
in \emph{International Conference on Machine Learning (ICML)},
2021.
[Online]. Available:
\url{https://arxiv.org/abs/2103.00020}

\bibitem{ranjan2021fsc147}
V.~Ranjan, U.~Sharma, T.~Nguyen, and M.~Hoai,
``Learning to count everything (FSC-147),''
in \emph{IEEE/CVF Conference on Computer Vision and Pattern
Recognition (CVPR)}, 2021.
[Online]. Available:
\url{https://arxiv.org/abs/2104.08391}

\bibitem{gemini}
Google DeepMind, ``Gemini API documentation,'' 2025.
[Online]. Available:
\url{https://ai.google.dev/gemini-api/docs}

\bibitem{streamlit}
Streamlit, ``The fastest way to build and share data apps,'' 2026.
[Online]. Available: \url{https://streamlit.io}

\bibitem{mae}
K.~He, X.~Chen, S.~Xie, Y.~Li, P.~Doll{\'a}r, and R.~Girshick,
``Masked autoencoders are scalable vision learners,''
in \emph{IEEE/CVF Conference on Computer Vision and Pattern
Recognition (CVPR)}, 2022.
[Online]. Available:
\url{https://arxiv.org/abs/2111.06377}

\bibitem{famnet}
M.~Ranasinghe, S.~Herath, H.~Bilen, and F.~Zheng,
``Learning to count without annotations,''
in \emph{IEEE/CVF Conference on Computer Vision and Pattern
Recognition (CVPR)}, 2022.
[Online]. Available:
\url{https://arxiv.org/abs/2307.08727}

\bibitem{bmnet}
M.~Shi, H.~Lu, C.~Feng, C.~Liu, and Z.~Cao,
``Represent, compare, and learn: A similarity-aware framework
for class-agnostic counting,''
in \emph{IEEE/CVF Conference on Computer Vision and Pattern
Recognition (CVPR)}, 2022.
[Online]. Available:
\url{https://arxiv.org/abs/2203.08354}

\bibitem{simclr}
T.~Chen, S.~Kornblith, M.~Norouzi, and G.~Hinton,
``A simple framework for contrastive learning of visual
representations (SimCLR),''
in \emph{International Conference on Machine Learning (ICML)}, 2020.
[Online]. Available:
\url{https://arxiv.org/abs/2002.05709}

\bibitem{zeroshot_counting_xu}
V.~Xu, L.~Zhu, and Y.~Yang,
``Zero-shot object counting,''
in \emph{IEEE/CVF Conference on Computer Vision and Pattern
Recognition (CVPR)}, 2023.
[Online]. Available:
\url{https://arxiv.org/abs/2303.02001}

\bibitem{coop}
K.~Zhou, J.~Yang, C.~C.~Loy, and Z.~Liu,
``Learning to prompt for vision-language models (CoOp),''
\emph{International Journal of Computer Vision (IJCV)}, vol.~130,
no.~9, pp.~2337--2348, 2022.
[Online]. Available:
\url{https://arxiv.org/abs/2109.01134}

\bibitem{ovdet_survey}
Y.~Zang, W.~Li, J.~Han, K.~Zhou, and C.~C.~Loy,
``Open-vocabulary object detection: A survey,''
arXiv preprint arXiv:2312.15850, 2023.
[Online]. Available:
\url{https://arxiv.org/abs/2312.15850}

\end{thebibliography}

\end{document}
